% ======================================================================
%  MODERN AI FOUNDATIONS -- A COMPLETE OVERVIEW
%  A self-contained course / study guide in LaTeX.
%
%  HOW TO COMPILE (pdfLaTeX + Biber):
%      1) pdflatex ai-foundations.tex
%      2) biber    ai-foundations
%      3) pdflatex ai-foundations.tex
%      4) pdflatex ai-foundations.tex
%
%  Requirements: TeX Live 2019+ (or MiKTeX) with the packages
%  tcolorbox, fontawesome5, pgfplots, FiraSans, FiraMono, biblatex
%  (biber), listings, titlesec, booktabs, enumitem, fancyhdr, csquotes.
%
%  If the Fira fonts are unavailable, replace the two font packages
%  below with \usepackage{lmodern}.
% ======================================================================
\begin{filecontents*}[overwrite]{ai-foundations.bib}
	@article{turing1950,
		author  = {Turing, Alan M.},
		title   = {Computing Machinery and Intelligence},
		journal = {Mind},
		volume  = {59}, number = {236}, pages = {433--460}, year = {1950}}
	@article{searle1980,
		author  = {Searle, John R.},
		title   = {Minds, Brains, and Programs},
		journal = {Behavioral and Brain Sciences},
		volume  = {3}, number = {3}, pages = {417--424}, year = {1980}}
	@book{russell2021,
		author  = {Russell, Stuart J. and Norvig, Peter},
		title   = {Artificial Intelligence: A Modern Approach},
		edition = {4}, publisher = {Pearson}, address = {Hoboken, NJ}, year = {2021}}
	@book{mitchell1997,
		author  = {Mitchell, Tom M.},
		title   = {Machine Learning},
		publisher = {McGraw-Hill}, address = {New York}, year = {1997}}
	@book{goodfellow2016,
		author  = {Goodfellow, Ian and Bengio, Yoshua and Courville, Aaron},
		title   = {Deep Learning},
		publisher = {MIT Press}, address = {Cambridge, MA}, year = {2016}}
	@book{bishop2006,
		author  = {Bishop, Christopher M.},
		title   = {Pattern Recognition and Machine Learning},
		publisher = {Springer}, address = {New York}, year = {2006}}
	@book{bishop2024,
		author  = {Bishop, Christopher M. and Bishop, Hugh},
		title   = {Deep Learning: Foundations and Concepts},
		publisher = {Springer}, address = {Cham}, year = {2024}}
	@book{hastie2009,
		author  = {Hastie, Trevor and Tibshirani, Robert and Friedman, Jerome},
		title   = {The Elements of Statistical Learning: Data Mining, Inference, and Prediction},
		edition = {2}, publisher = {Springer}, address = {New York}, year = {2009}}
	@book{strang2019,
		author  = {Strang, Gilbert},
		title   = {Linear Algebra and Learning from Data},
		publisher = {Wellesley--Cambridge Press}, address = {Wellesley, MA}, year = {2019}}
	@book{wasserman2004,
		author  = {Wasserman, Larry},
		title   = {All of Statistics: A Concise Course in Statistical Inference},
		publisher = {Springer}, address = {New York}, year = {2004}}
	@book{boyd2004,
		author  = {Boyd, Stephen and Vandenberghe, Lieven},
		title   = {Convex Optimization},
		publisher = {Cambridge University Press}, address = {Cambridge}, year = {2004}}
	@book{christian2020,
		author  = {Christian, Brian},
		title   = {The Alignment Problem: Machine Learning and Human Values},
		publisher = {W.\,W.\,Norton \& Company}, address = {New York}, year = {2020}}
	@book{vanderplas2016,
		author  = {VanderPlas, Jake},
		title   = {Python Data Science Handbook: Essential Tools for Working with Data},
		publisher = {O'Reilly Media}, address = {Sebastopol, CA}, year = {2016}}
	@article{lecun2015,
		author  = {LeCun, Yann and Bengio, Yoshua and Hinton, Geoffrey},
		title   = {Deep Learning},
		journal = {Nature},
		volume  = {521}, number = {7553}, pages = {436--444}, year = {2015}}
	@inproceedings{krizhevsky2012,
		author  = {Krizhevsky, Alex and Sutskever, Ilya and Hinton, Geoffrey E.},
		title   = {ImageNet Classification with Deep Convolutional Neural Networks},
		booktitle = {Advances in Neural Information Processing Systems 25},
		pages   = {1097--1105}, year = {2012}}
	@article{silver2016,
		author  = {Silver, David and Huang, Aja and Maddison, Chris J. and others},
		title   = {Mastering the Game of {Go} with Deep Neural Networks and Tree Search},
		journal = {Nature},
		volume  = {529}, number = {7587}, pages = {484--489}, year = {2016}}
	@inproceedings{jouppi2017,
		author  = {Jouppi, Norman P. and Young, Cliff and Patil, Nishant and others},
		title   = {In-Datacenter Performance Analysis of a Tensor Processing Unit},
		booktitle = {Proceedings of the 44th Annual International Symposium on Computer Architecture (ISCA)},
		pages   = {1--12}, year = {2017}}
	@article{harris2020,
		author  = {Harris, Charles R. and Millman, K. Jarrod and van der Walt, St{\'e}fan J. and others},
		title   = {Array Programming with {NumPy}},
		journal = {Nature},
		volume  = {585}, pages = {357--362}, year = {2020}}
	@inproceedings{mckinney2010,
		author  = {McKinney, Wes},
		title   = {Data Structures for Statistical Computing in {Python}},
		booktitle = {Proceedings of the 9th Python in Science Conference (SciPy)},
		pages   = {56--61}, year = {2010}}
	@article{pedregosa2011,
		author  = {Pedregosa, Fabian and Varoquaux, Ga{\"e}l and Gramfort, Alexandre and others},
		title   = {Scikit-learn: Machine Learning in {Python}},
		journal = {Journal of Machine Learning Research},
		volume  = {12}, pages = {2825--2830}, year = {2011}}
	@article{rosenblatt1958,
		author  = {Rosenblatt, Frank},
		title   = {The Perceptron: A Probabilistic Model for Information Storage and Organization in the Brain},
		journal = {Psychological Review},
		volume  = {65}, number = {6}, pages = {386--408}, year = {1958}}
	@book{minsky1969,
		author  = {Minsky, Marvin and Papert, Seymour},
		title   = {Perceptrons: An Introduction to Computational Geometry},
		publisher = {MIT Press}, address = {Cambridge, MA}, year = {1969}}
	@article{rumelhart1986,
		author  = {Rumelhart, David E. and Hinton, Geoffrey E. and Williams, Ronald J.},
		title   = {Learning Representations by Back-Propagating Errors},
		journal = {Nature},
		volume  = {323}, number = {6088}, pages = {533--536}, year = {1986}}
	@article{hornik1989,
		author  = {Hornik, Kurt and Stinchcombe, Maxwell and White, Halbert},
		title   = {Multilayer Feedforward Networks Are Universal Approximators},
		journal = {Neural Networks},
		volume  = {2}, number = {5}, pages = {359--366}, year = {1989}}
	@inproceedings{nair2010,
		author  = {Nair, Vinod and Hinton, Geoffrey E.},
		title   = {Rectified Linear Units Improve Restricted {Boltzmann} Machines},
		booktitle = {Proceedings of the 27th International Conference on Machine Learning (ICML)},
		pages   = {807--814}, year = {2010}}
	@article{lecun1998,
		author  = {LeCun, Yann and Bottou, L{\'e}on and Bengio, Yoshua and Haffner, Patrick},
		title   = {Gradient-Based Learning Applied to Document Recognition},
		journal = {Proceedings of the IEEE},
		volume  = {86}, number = {11}, pages = {2278--2324}, year = {1998}}
	@article{hochreiter1997,
		author  = {Hochreiter, Sepp and Schmidhuber, J{\"u}rgen},
		title   = {Long Short-Term Memory},
		journal = {Neural Computation},
		volume  = {9}, number = {8}, pages = {1735--1780}, year = {1997}}
	@inproceedings{vaswani2017,
		author  = {Vaswani, Ashish and Shazeer, Noam and Parmar, Niki and Uszkoreit, Jakob
			and Jones, Llion and Gomez, Aidan N. and Kaiser, {\L}ukasz and Polosukhin, Illia},
		title   = {Attention Is All You Need},
		booktitle = {Advances in Neural Information Processing Systems 30},
		pages   = {5998--6008}, year = {2017}}
	@inproceedings{mikolov2013,
		author  = {Mikolov, Tom{\'a}{\v{s}} and Sutskever, Ilya and Chen, Kai
			and Corrado, Greg S. and Dean, Jeffrey},
		title   = {Distributed Representations of Words and Phrases and Their Compositionality},
		booktitle = {Advances in Neural Information Processing Systems 26},
		pages   = {3111--3119}, year = {2013}}
	@inproceedings{kingma2014,
		author  = {Kingma, Diederik P. and Welling, Max},
		title   = {Auto-Encoding Variational {Bayes}},
		booktitle = {International Conference on Learning Representations (ICLR)},
		pages   = {1--9}, year = {2014}}
	@inproceedings{goodfellow2014,
		author  = {Goodfellow, Ian and Pouget-Abadie, Jean and Mirza, Mehdi and Xu, Bing
			and Warde-Farley, David and Ozair, Sherjil and Courville, Aaron and Bengio, Yoshua},
		title   = {Generative Adversarial Nets},
		booktitle = {Advances in Neural Information Processing Systems 27},
		pages   = {2672--2680}, year = {2014}}
	@inproceedings{brown2020,
		author  = {Brown, Tom B. and Mann, Benjamin and Ryder, Nick and others},
		title   = {Language Models Are Few-Shot Learners},
		booktitle = {Advances in Neural Information Processing Systems 33},
		pages   = {1877--1901}, year = {2020}}
	@article{kaplan2020,
		author  = {Kaplan, Jared and McCandlish, Sam and Henighan, Tom and others},
		title   = {Scaling Laws for Neural Language Models},
		journal = {arXiv preprint arXiv:2001.08361},
		year = {2020}}
	@inproceedings{ouyang2022,
		author  = {Ouyang, Long and Wu, Jeffrey and Jiang, Xu and others},
		title   = {Training Language Models to Follow Instructions with Human Feedback},
		booktitle = {Advances in Neural Information Processing Systems 35},
		pages   = {27730--27744}, year = {2022}}
	@inproceedings{christiano2017,
		author  = {Christiano, Paul F. and Leike, Jan and Brown, Tom B.
			and Martic, Miljan and Legg, Shane and Amodei, Dario},
		title   = {Deep Reinforcement Learning from Human Preferences},
		booktitle = {Advances in Neural Information Processing Systems 30},
		pages   = {4299--4307}, year = {2017}}
	@inproceedings{bender2021,
		author  = {Bender, Emily M. and Gebru, Timnit and McMillan-Major, Angelina
			and Shmitchell, Shmargaret},
		title   = {On the Dangers of Stochastic Parrots: Can Language Models Be Too Big?},
		booktitle = {Proceedings of the 2021 ACM Conference on Fairness, Accountability, and Transparency (FAccT)},
		pages   = {610--623}, year = {2021}}
	@inproceedings{buolamwini2018,
		author  = {Buolamwini, Joy and Gebru, Timnit},
		title   = {Gender Shades: Intersectional Accuracy Disparities in Commercial Gender Classification},
		booktitle = {Proceedings of the 1st Conference on Fairness, Accountability and Transparency},
		pages   = {77--91}, year = {2018}}
	@article{amodei2016,
		author  = {Amodei, Dario and Olah, Chris and Steinhardt, Jacob and others},
		title   = {Concrete Problems in {AI} Safety},
		journal = {arXiv preprint arXiv:1606.06565},
		year = {2016}}
	@inproceedings{strubell2019,
		author  = {Strubell, Emma and Ganesh, Ananya and McCallum, Andrew},
		title   = {Energy and Policy Considerations for Deep Learning in {NLP}},
		booktitle = {Proceedings of the 57th Annual Meeting of the Association for Computational Linguistics (ACL)},
		pages   = {3645--3650}, year = {2019}}
	@inproceedings{patterson2021,
		author  = {Patterson, David and Gonzalez, Joseph and Le, Quoc and others},
		title   = {Carbon Emissions and Large Neural Network Training},
		booktitle = {Advances in Neural Information Processing Systems 34},
		pages   = {13470--13481}, year = {2021}}
	@article{cortes1995,
		author  = {Cortes, Corinna and Vapnik, Vladimir},
		title   = {Support-Vector Networks},
		journal = {Machine Learning},
		volume  = {20}, number = {3}, pages = {273--297}, year = {1995}}
	@article{breiman2001,
		author  = {Breiman, Leo},
		title   = {Random Forests},
		journal = {Machine Learning},
		volume  = {45}, number = {1}, pages = {5--32}, year = {2001}}
	@inproceedings{chen2016,
		author  = {Chen, Tianqi and Guestrin, Carlos},
		title   = {{XGBoost}: A Scalable Tree Boosting System},
		booktitle = {Proceedings of the 22nd ACM SIGKDD International Conference on Knowledge Discovery and Data Mining},
		pages   = {785--794}, year = {2016}}
	@inproceedings{macqueen1967,
		author  = {MacQueen, James},
		title   = {Some Methods for Classification and Analysis of Multivariate Observations},
		booktitle = {Proceedings of the Fifth Berkeley Symposium on Mathematical Statistics and Probability},
		volume  = {1}, pages = {281--297}, year = {1967}}
	@inproceedings{liu2008,
		author  = {Liu, Fei Tony and Ting, Kai Ming and Zhou, Zhi-Hua},
		title   = {Isolation Forest},
		booktitle = {Proceedings of the 8th IEEE International Conference on Data Mining (ICDM)},
		pages   = {413--422}, year = {2008}}
	@article{fawcett2006,
		author  = {Fawcett, Tom},
		title   = {An Introduction to {ROC} Analysis},
		journal = {Pattern Recognition Letters},
		volume  = {27}, number = {8}, pages = {861--874}, year = {2006}}
	@book{zhang2023,
		author  = {Zhang, Aston and Lipton, Zachary C. and Li, Mu and Smola, Alexander J.},
		title   = {Dive into Deep Learning},
		publisher = {Cambridge University Press}, address = {Cambridge}, year = {2023}}
	@online{fastai,
		author  = {Howard, Jeremy and others},
		title   = {Practical Deep Learning for Coders},
		organization = {fast.ai},
		url = {https://course.fast.ai},
		urldate = {2024-12-01}}
	@online{huggingface,
		author  = {{Hugging Face}},
		title   = {Transformers Documentation},
		url = {https://huggingface.co/docs/transformers},
		urldate = {2024-12-01}}
	@inproceedings{kingma2015,
		author  = {Kingma, Diederik P. and Ba, Jimmy},
		title   = {Adam: A Method for Stochastic Optimization},
		booktitle = {International Conference on Learning Representations (ICLR)},
		pages   = {1--15}, year = {2015}}
	@inproceedings{he2016,
		author  = {He, Kaiming and Zhang, Xiangyu and Ren, Shaoqing and Sun, Jian},
		title   = {Deep Residual Learning for Image Recognition},
		booktitle = {Proceedings of the IEEE Conference on Computer Vision and Pattern Recognition (CVPR)},
		pages   = {770--778}, year = {2016}}
	@article{su2024,
		author  = {Su, Jianlin and Ahmed, Murtadha and Lu, Yu and Pan, Shengfeng and Bo, Wen and Liu, Yunfeng},
		title   = {RoFormer: Enhanced Transformer with Rotary Position Embedding},
		journal = {Neurocomputing},
		volume  = {568}, pages = {127063}, year = {2024}}
	@article{rousseeuw1987,
		author  = {Rousseeuw, Peter J.},
		title   = {Silhouettes: A Graphical Aid to the Interpretation and Validation of Cluster Analysis},
		journal = {Journal of Computational and Applied Mathematics},
		volume  = {20}, pages = {53--65}, year = {1987}}
	@inproceedings{wolf2020,
		author  = {Wolf, Thomas and Debut, Lysandre and Sanh, Victor and Chaumond, Julien and Delangue, Clement and Moi, Anthony and Cistac, Pierric and Rault, Timoth{\'e}e and Louf, R{\'e}mi and Funtowicz, Morgan and Davison, Joe and Shleifer, Sam and von Platen, Patrick and Ma, Clara and Jernite, Yacine and Plu, Julien and Xu, Canwen and Le Scao, Teven and Gugger, Sylvain and Drame, Mariama and Lhoest, Quentin and Rush, Alexander M.},
		title   = {Transformers: State-of-the-Art Natural Language Processing},
		booktitle = {Proceedings of the 2020 Conference on Empirical Methods in Natural Language Processing: System Demonstrations},
		pages   = {38--45}, year = {2020}}
	@article{hoffmann2022,
		author  = {Hoffmann, Jordan and Borgeaud, Sebastian and Mensch, Arthur and Buchatskaya, Elena and Cai, Trevor and Rutherford, Eliza and de Las Casas, Diego and Hendricks, Lisa Anne and Welbl, Johannes and Clark, Aidan and Hennigan, Tom and Noland, Eric and Millican, Katie and Driessche, George van den and Damoc, Bogdan and Guy, Aurelia and Osindero, Simon and Simonyan, Karen and Elsen, Erich and Rae, Jack W. and Vinyals, Oriol and Sifre, Laurent},
		title   = {Training Compute-Optimal Large Language Models},
		journal = {arXiv preprint arXiv:2203.15556},
		year    = {2022}}
	@inproceedings{ho2020,
		author  = {Ho, Jonathan and Jain, Ajay and Abbeel, Pieter},
		title   = {Denoising Diffusion Probabilistic Models},
		booktitle = {Advances in Neural Information Processing Systems 33},
		pages   = {6840--6851}, year = {2020}}
	@book{sutton2018,
		author    = {Sutton, Richard S. and Barto, Andrew G.},
		title     = {Reinforcement Learning: An Introduction},
		edition   = {2},
		publisher = {MIT Press},
		address   = {Cambridge, MA},
		year      = {2018}}
	@article{watkins1992,
		author  = {Watkins, Christopher J. C. H. and Dayan, Peter},
		title   = {Q-Learning},
		journal = {Machine Learning},
		volume  = {8}, number = {3-4}, pages = {279--292}, year = {1992}}
	@article{mnih2015,
		author  = {Mnih, Volodymyr and Kavukcuoglu, Koray and Silver, David and others},
		title   = {Human-level Control Through Deep Reinforcement Learning},
		journal = {Nature},
		volume  = {518}, number = {7540}, pages = {529--533}, year = {2015}}
	@inproceedings{schulman2017,
		author  = {Schulman, John and Wolski, Filip and Dhariwal, Prafulla and Radford, Alec and Klimov, Oleg},
		title   = {Proximal Policy Optimization Algorithms},
		journal = {arXiv preprint arXiv:1707.06347},
		year    = {2017}}
	@inproceedings{rafailov2023,
		author  = {Rafailov, Rafael and Sharma, Archit and Mitchell, Eric and Ermon, Stefano and Manning, Christopher D. and Finn, Chelsea},
		title   = {Direct Preference Optimization: Your Language Model is Secretly a Reward Model},
		booktitle = {Advances in Neural Information Processing Systems 36},
		year    = {2023}}
	@inproceedings{hu2021,
		author  = {Hu, Edward J. and Shen, Yelong and Wallis, Phillip and Allen-Zhu, Zeyuan and Li, Yuanzhi and Wang, Shean and Wang, Lu and Chen, Weizhu},
		title   = {{LoRA}: Low-Rank Adaptation of Large Language Models},
		booktitle = {International Conference on Learning Representations (ICLR)},
		year    = {2022}}
	@inproceedings{dettmers2024,
		author  = {Dettmers, Tim and Pagnoni, Artidoro and Holtzman, Ari and Zettlemoyer, Luke},
		title   = {{QLoRA}: Efficient Finetuning of Quantized {LLMs}},
		booktitle = {Advances in Neural Information Processing Systems 36},
		year    = {2024}}
	@article{dubey2024,
		author  = {Dubey, Abhimanyu and Jauhri, Abhinav and Pandey, Abhinav and others},
		title   = {The {Llama} 3 Herd of Models},
		journal = {arXiv preprint arXiv:2407.21783},
		year    = {2024}}
	@inproceedings{lewis2020,
		author  = {Lewis, Patrick and Perez, Ethan and Piktus, Aleksandra and Petroni, Fabio and Karpukhin, Vladimir and Goyal, Naman and Kuttler, Heinrich and Lewis, Mike and Yih, Wen-tau and Rockt{\"a}schel, Tim and Riedel, Sebastian and Kiela, Douwe},
		title   = {Retrieval-Augmented Generation for Knowledge-Intensive {NLP} Tasks},
		booktitle = {Advances in Neural Information Processing Systems 33},
		pages   = {9459--9474}, year = {2020}}
	@inproceedings{ester1996,
		author  = {Ester, Martin and Kriegel, Hans-Peter and Sander, J{\"o}rg and Xu, Xiaowei},
		title   = {A Density-Based Algorithm for Discovering Clusters in Large Spatial Databases with Noise},
		booktitle = {Proceedings of the 2nd International Conference on Knowledge Discovery and Data Mining (KDD)},
		pages   = {226--231}, year = {1996}}
	@article{vanDerMaaten2008,
		author  = {van der Maaten, Laurens and Hinton, Geoffrey},
		title   = {Visualizing Data using {t-SNE}},
		journal = {Journal of Machine Learning Research},
		volume  = {9}, pages = {2579--2605}, year = {2008}}
	@article{mcinnes2018,
		author  = {McInnes, Leland and Healy, John and Melville, James},
		title   = {{UMAP}: Uniform Manifold Approximation and Projection for Dimension Reduction},
		journal = {arXiv preprint arXiv:1802.03426},
		year    = {2018}}
	@inproceedings{glorot2010,
		author  = {Glorot, Xavier and Bengio, Yoshua},
		title   = {Understanding the Difficulty of Training Deep Feedforward Neural Networks},
		booktitle = {Proceedings of the 13th International Conference on Artificial Intelligence and Statistics (AISTATS)},
		pages   = {249--256}, year = {2010}}
	@inproceedings{he2015,
		author  = {He, Kaiming and Zhang, Xiangyu and Ren, Shaoqing and Sun, Jian},
		title   = {Delving Deep into Rectifiers: Surpassing Human-Level Performance on {ImageNet} Classification},
		booktitle = {Proceedings of the IEEE International Conference on Computer Vision (ICCV)},
		pages   = {1026--1034}, year = {2015}}
	@inproceedings{loshchilov2019,
		author  = {Loshchilov, Ilya and Hutter, Frank},
		title   = {Decoupled Weight Decay Regularization},
		booktitle = {International Conference on Learning Representations (ICLR)},
		year    = {2019}}
	@article{bengio2013,
		author  = {Bengio, Yoshua and Courville, Aaron and Vincent, Pascal},
		title   = {Representation Learning: A Review and New Perspectives},
		journal = {IEEE Transactions on Pattern Analysis and Machine Intelligence},
		volume  = {35}, number = {8}, pages = {1798--1828}, year = {2013}}
	@inproceedings{vincent2008,
		author  = {Vincent, Pascal and Larochelle, Hugo and Bengio, Yoshua and Manzagol, Pierre-Antoine},
		title   = {Extracting and Composing Robust Features with Denoising Autoencoders},
		booktitle = {Proceedings of the 25th International Conference on Machine Learning (ICML)},
		pages   = {1096--1103}, year = {2008}}
	@article{oord2018,
		author  = {van den Oord, Aaron and Li, Yazhe and Vinyals, Oriol},
		title   = {Representation Learning with Contrastive Predictive Coding},
		journal = {arXiv preprint arXiv:1807.03748},
		year    = {2018}}
	@inproceedings{chen2020,
		author  = {Chen, Ting and Kornblith, Simon and Norouzi, Mohammad and Hinton, Geoffrey},
		title   = {A Simple Framework for Contrastive Learning of Visual Representations},
		booktitle = {Proceedings of the 37th International Conference on Machine Learning (ICML)},
		pages   = {1597--1607}, year = {2020}}
	@inproceedings{devlin2019,
		author  = {Devlin, Jacob and Chang, Ming-Wei and Lee, Kenton and Toutanova, Kristina},
		title   = {{BERT}: Pre-training of Deep Bidirectional Transformers for Language Understanding},
		booktitle = {Proceedings of the 2019 Conference of the North American Chapter of the Association for Computational Linguistics (NAACL)},
		pages   = {4171--4186}, year = {2019}}
	@article{radford2018,
		author  = {Radford, Alec and Narasimhan, Karthik and Salimans, Tim and Sutskever, Ilya},
		title   = {Improving Language Understanding by Generative Pre-Training},
		journal = {OpenAI Blog},
		year    = {2018}}
	@inproceedings{he2022,
		author  = {He, Kaiming and Chen, Xinlei and Xie, Saining and Li, Yanghao and Doll{\'a}r, Piotr and Girshick, Ross},
		title   = {Masked Autoencoders Are Scalable Vision Learners},
		booktitle = {Proceedings of the IEEE/CVF Conference on Computer Vision and Pattern Recognition (CVPR)},
		pages   = {16000--16009}, year = {2022}}
	@inproceedings{dosovitskiy2021,
		author  = {Dosovitskiy, Alexey and Beyer, Lucas and Kolesnikov, Alexander and Weissenborn, Dirk and Zhai, Xiaohua and Unterthiner, Thomas and Dehghani, Mostafa and Minderer, Matthias and Heigold, Georg and Gelly, Sylvain and Uszkoreit, Jakob and Houlsby, Neil},
		title   = {An Image is Worth 16x16 Words: Transformers for Image Recognition at Scale},
		booktitle = {International Conference on Learning Representations (ICLR)},
		year    = {2021}}
\end{filecontents*}

\documentclass[11pt,a4paper]{article}

% ---------------------------------------------------------------------
% Fonts & typography
% ---------------------------------------------------------------------
\usepackage[T1]{fontenc}
\usepackage[utf8]{inputenc}
\usepackage[sfdefault]{FiraSans}   % fallback: \usepackage{lmodern}
\usepackage[scaled=0.86]{FiraMono} % fallback: remove this line
\usepackage{microtype}

% ---------------------------------------------------------------------
% Layout
% ---------------------------------------------------------------------
\usepackage[a4paper,margin=2.35cm,headheight=15pt]{geometry}
\setlength{\parindent}{0pt}
\setlength{\parskip}{4.5pt plus 1pt}
\raggedbottom
\emergencystretch=1.5em

% ---------------------------------------------------------------------
% Colour palette & tables
% ---------------------------------------------------------------------
\usepackage[table]{xcolor}
\definecolor{primary}{RGB}{38,60,114}   % deep indigo
\definecolor{secondary}{RGB}{0,150,136} % teal
\definecolor{accent}{RGB}{230,126,34}   % amber
\definecolor{dpurple}{RGB}{108,52,158}  % deep purple
\definecolor{textgray}{RGB}{75,80,90}

\usepackage{booktabs,array}
\newcolumntype{L}[1]{>{\raggedright\arraybackslash}p{#1}}

\usepackage{enumitem}
\setlist[itemize]{topsep=2.5pt,itemsep=2.5pt,parsep=0pt,leftmargin=1.5em}
\setlist[enumerate]{topsep=2.5pt,itemsep=2.5pt,parsep=0pt,leftmargin=1.7em}
\setlist[description]{topsep=2.5pt,itemsep=3pt,parsep=0pt}

% ---------------------------------------------------------------------
% Math, graphics, diagrams
% ---------------------------------------------------------------------
\usepackage{amsmath,amssymb,mathtools}
\DeclareMathOperator{\softmax}{softmax}
\usepackage{graphicx}
\usepackage{tikz}
\usetikzlibrary{arrows.meta,positioning,calc}
\usepackage{pgfplots}
\pgfplotsset{compat=1.16}

% timeline milestone macro (used in Module 1)
\newcommand{\ms}[5]{% x, year, label, y-offset, anchor
	\draw[secondary, line width=1.1pt] (#1,-0.06) -- (#1,0.06);
	\node[align=center, font=\scriptsize\sffamily, anchor=#5] at (#1,#4)
	{\textbf{#2}\\#3};
}

% ---------------------------------------------------------------------
% Modern callout boxes
% ---------------------------------------------------------------------
\usepackage[most]{tcolorbox}
\usepackage{fontawesome5}

\newtcolorbox{modulecard}[2]{%
	enhanced, breakable,
	colback=#2!5, colframe=#2, boxrule=0.8pt, arc=2mm,
	left=10pt, right=10pt, top=8pt, bottom=8pt,
	toptitle=6pt, bottomtitle=6pt,
	fonttitle=\sffamily\bfseries\large,
	coltitle=white, colbacktitle=#2,
	before skip=12pt, after skip=10pt,
	title={#1}}

\newtcolorbox{keyconcept}[1][Key concept]{%
	enhanced, breakable,
	colback=primary!7, colframe=primary,
	boxrule=0pt, leftrule=3.5pt, arc=1.2mm,
	left=10pt, right=9pt, top=5pt, bottom=5pt,
	toptitle=3pt, bottomtitle=2pt,
	fonttitle=\sffamily\bfseries, coltitle=primary, colbacktitle=primary!7,
	title={\faLightbulb\ #1},
	before skip=10pt, after skip=12pt}

\newtcolorbox{practicebox}[1][Hands-on with Python]{%
	enhanced,
	colback=secondary!7, colframe=secondary,
	boxrule=0pt, leftrule=3.5pt, arc=1.2mm,
	left=10pt, right=9pt, top=5pt, bottom=5pt,
	toptitle=3pt, bottomtitle=2pt,
	fonttitle=\sffamily\bfseries, coltitle=secondary, colbacktitle=secondary!7,
	title={\faPython\ #1},
	before skip=10pt, after skip=12pt}

\newtcolorbox{ethicsbox}[1][Ethics \& safety spotlight]{%
	enhanced, breakable,
	colback=accent!8, colframe=accent,
	boxrule=0pt, leftrule=3.5pt, arc=1.2mm,
	left=10pt, right=9pt, top=5pt, bottom=5pt,
	toptitle=3pt, bottomtitle=2pt,
	fonttitle=\sffamily\bfseries, coltitle=accent!60!black, colbacktitle=accent!8,
	title={\faBalanceScale\ #1},
	before skip=10pt, after skip=12pt}

\newtcolorbox{checkpoint}[1][Check your understanding]{%
	enhanced, breakable,
	colback=dpurple!6, colframe=dpurple!70,
	boxrule=0.7pt, arc=1.2mm,
	left=10pt, right=9pt, top=5pt, bottom=5pt,
	toptitle=3pt, bottomtitle=2pt,
	fonttitle=\sffamily\bfseries, coltitle=dpurple, colbacktitle=dpurple!6,
	title={\faQuestionCircle\ #1},
	before skip=10pt, after skip=12pt}

\newtcolorbox{sourcesbox}[1][Sources \& further reading]{%
	enhanced, breakable,
	colback=black!4, colframe=black!35,
	boxrule=0.5pt, arc=1mm,
	left=9pt, right=9pt, top=5pt, bottom=5pt,
	toptitle=3pt, bottomtitle=2pt,
	fonttitle=\sffamily\bfseries, coltitle=black!70, colbacktitle=black!4,
	title={\faBookOpen\ #1},
	before skip=10pt, after skip=14pt}

% ---------------------------------------------------------------------
% Code listings
% ---------------------------------------------------------------------
\usepackage{listings}
\lstset{
	language=Python,
	basicstyle=\ttfamily\footnotesize,
	keywordstyle=\color{secondary}\bfseries,
	keywordstyle=[2]\color{primary},
	morekeywords=[2]{np,pd,plt,sklearn,rng},
	stringstyle=\color{accent!85!black},
	commentstyle=\color{textgray}\itshape,
	showstringspaces=false,
	breaklines=true,
	breakatwhitespace=true,
	columns=flexible,
	keepspaces=true,
	frame=leftline,
	framerule=1.2pt,
	rulecolor=\color{secondary!55},
	xleftmargin=0.5em,
	aboveskip=0.7em, belowskip=0.7em}

% ---------------------------------------------------------------------
% Headings, headers, captions
% ---------------------------------------------------------------------
\usepackage{titlesec}
\titleformat{\section}
{\Large\bfseries\sffamily\color{primary}}{\thesection}{0.7em}{}
[{\color{secondary}\titlerule[1.1pt]}]
\titleformat{\subsection}
{\large\bfseries\sffamily\color{primary!85!black}}{\thesubsection}{0.6em}{}
\titleformat{\subsubsection}
{\normalsize\bfseries\sffamily\color{secondary}}{\thesubsubsection}{0.5em}{}
\titlespacing*{\section}{0pt}{3.2ex plus 1ex}{2.2ex}
\titlespacing*{\subsection}{0pt}{2.6ex plus 0.8ex}{1.4ex}

\usepackage{fancyhdr}
\pagestyle{fancy}
\fancyhf{}
\fancyhead[L]{\small\sffamily\color{primary}\faBrain\ Modern AI Foundations}
\fancyhead[R]{\small\sffamily\color{textgray}\nouppercase\rightmark}
\fancyfoot[C]{\small\sffamily\color{textgray}\thepage}
\renewcommand{\headrulewidth}{0pt}
\renewcommand{\headrule}{{\color{secondary!60}\hrule width\headwidth height 0.8pt}}
\renewcommand{\sectionmark}[1]{\markright{\thesection.\ #1}}

\usepackage{caption}
\captionsetup{font=small, labelfont={bf,color=primary}, labelsep=period}
\captionsetup[table]{skip=5pt}

% ---------------------------------------------------------------------
% Bibliography (biblatex + biber)
% ---------------------------------------------------------------------
\usepackage{csquotes}
\usepackage[backend=biber, style=authoryear, giveninits=true,
maxcitenames=2, maxbibnames=6, uniquename=false,
uniquelist=false, dashed=false]{biblatex}
\addbibresource{ai-foundations.bib}
\renewcommand{\bibfont}{\small}
\setlength{\bibitemsep}{4pt}

% ---------------------------------------------------------------------
% Hyperlinks (load last)
% ---------------------------------------------------------------------
\usepackage[colorlinks=true, linkcolor=primary, citecolor=secondary,
urlcolor=secondary, bookmarksnumbered=true,
pdftitle={Modern AI Foundations: A Complete Overview},
pdfauthor={Course Design Team},
pdfsubject={Artificial Intelligence, Machine Learning, Deep Learning},
pdfkeywords={AI, machine learning, deep learning, transformers, LLM}]{hyperref}
\urlstyle{same}

% =====================================================================
\begin{document}
	% =====================================================================
	
	% ---------------------------------------------------------------------
	% Title page
	% ---------------------------------------------------------------------
	\hypersetup{pageanchor=false}
	\begin{titlepage}
		\centering
		\begin{tikzpicture}[remember picture, overlay]
			\fill[primary]   (current page.north west) rectangle ([yshift=-4.7cm]current page.north east);
			\fill[secondary] ([yshift=-4.7cm]current page.north west) rectangle ([yshift=-5.0cm]current page.north east);
			\fill[primary]   (current page.south west) rectangle ([yshift=1.9cm]current page.south east);
			\fill[secondary] ([yshift=1.9cm]current page.south west) rectangle ([yshift=2.15cm]current page.south east);
			\node[align=center, text=white, font=\sffamily\bfseries\fontsize{28}{33}\selectfont]
			at ([yshift=-2.1cm]current page.north) {Modern AI Foundations};
			\node[align=center, text=white!85!primary, font=\sffamily\Large, text width=14.5cm]
			at ([yshift=-3.7cm]current page.north)
			{An Accelerated Introduction to AI: From Classical Machine Learning\\to Large Language Models};
		\end{tikzpicture}
		\vspace*{5.4cm}
		
		{\color{secondary}\large\faGlobe\hspace{1.2em}\faChartLine\hspace{1.2em}
			\faCogs\hspace{1.2em}\faRobot\hspace{1.2em}\faBrain}\\[1.0em]
		{\color{primary}\rule{0.55\textwidth}{1pt}}\\[1.6em]
		
		\begin{minipage}{0.86\textwidth}
			\centering\small\itshape
			An accelerated, code-first guide to the ideas behind modern artificial
			intelligence: the mathematics that powers it, the classical
			machine-learning engine at its heart, and the deep-learning revolution
			that brought us transformers and large language models --- with curated
			sources for every step of the journey.
		\end{minipage}\\[2.4em]
		
		{\sffamily\color{primary}\bfseries
			Prerequisites: Intermediate Python \textbullet{}
			Basic NumPy \& Pandas \textbullet{} Calculus \& Linear Algebra Basics}\\[0.9em]
		
		\vfill
		{\sffamily\color{textgray}\large Course Design Team}\\[0.3em]
		{\sffamily\color{textgray}\small Version 1.1 \textbullet{} September 2025 \textbullet{} License: CC BY 4.0}\\[0.3em]
		{\sffamily\color{secondary}\small Repository: \url{https://github.com/DanielBerns/AI-course}}
		\vspace*{2.2cm}
	\end{titlepage}
	\hypersetup{pageanchor=true}
	
	% ---------------------------------------------------------------------
	\clearpage
	\tableofcontents
	\clearpage
	
	% =====================================================================
	\section*{About this course}
	\markright{About this course}
	\phantomsection
	\addcontentsline{toc}{section}{About this course}
	% =====================================================================
	
	This document develops a five-module accelerated overview of modern artificial
	intelligence for technically capable students. It assumes basic background
	summarised in the prerequisite matrix on the next page: comfortable
	\textbf{Python}, basic \textbf{NumPy \& Pandas}, and \textbf{elementary
		calculus and linear algebra}. Everything else --- from eigenvectors to multi-head attention --- is built
	from scratch, in context, with citations to the primary literature and
	standard textbooks.
	
	The five modules are deliberately sequential, as the roadmap below shows:
	the landscape and its ethics (\emph{Module~1}), the mathematical
	foundations (\emph{Module~2}), the classical machine-learning engine that
	remains highly competitive for tabular workloads (\emph{Module~3}), neural networks and representation learning (\emph{Module~4}), and the foundation model and generative-AI frontier (\emph{Module~5}). All code snippets
	assume Python 3.9+ with \texttt{numpy}, \texttt{pandas},
	\texttt{scikit-learn}, and (for Modules~4 and~5) \texttt{transformers} installed.
	
	\begin{figure}[ht]
		\centering
		\begin{tikzpicture}[
			stage/.style={rounded corners=2mm, align=center, font=\sffamily\scriptsize,
				minimum height=1.05cm, inner sep=2pt, text width=2.5cm},
			>=Stealth]
			\node[stage, draw=primary,   fill=primary!8]   (m1) {\textbf{Module 1}\\[2pt] The AI landscape\\ context \& ethics};
			\node[stage, draw=secondary, fill=secondary!8, right=0.22cm of m1] (m2)
			{\textbf{Module 2}\\[2pt] Math \& data\\ language of ML};
			\node[stage, draw=accent,    fill=accent!8,    right=0.22cm of m2] (m3)
			{\textbf{Module 3}\\[2pt] Classical ML\\ the core engine};
			\node[stage, draw=dpurple,   fill=dpurple!8,   right=0.22cm of m3] (m4)
			{\textbf{Module 4}\\[2pt] Neural Networks\\ \& representation};
			\node[stage, draw=primary!80!black, fill=primary!15, right=0.22cm of m4] (m5)
			{\textbf{Module 5}\\[2pt] Foundation Models\\ \& generative AI};
			\draw[->, thick, primary!70]   (m1) -- (m2);
			\draw[->, thick, secondary!80] (m2) -- (m3);
			\draw[->, thick, accent!90]    (m3) -- (m4);
			\draw[->, thick, dpurple!90]   (m4) -- (m5);
		\end{tikzpicture}
		\caption{The learning path: each module builds directly on the one before it.}
		\label{fig:roadmap}
	\end{figure}
	
	Throughout the guide you will meet the following recurring callout boxes:
	
	\begin{table}[ht]
		\centering
		\small\sffamily
		\renewcommand{\arraystretch}{1.3}
		\begin{tabular}{@{}L{5.2cm} L{9.6cm}@{}}
			\toprule
			\rowcolor{primary}
			\textcolor{white}{\textbf{Box}} & \textcolor{white}{\textbf{What it gives you}} \\
			\midrule
			\textcolor{primary}{\faInfoCircle}\ Module --- at a glance & Focus, key concepts, and learning goals for each module. \\
			\textcolor{primary}{\faLightbulb}\ Key concept & Definitions and intuitions worth memorising. \\
			\textcolor{secondary}{\faPython}\ Hands-on with Python & Short, runnable code tied to the prerequisite skills. \\
			\textcolor{accent}{\faBalanceScale}\ Ethics \& safety spotlight & Responsible-AI notes and open questions. \\
			\textcolor{dpurple}{\faQuestionCircle}\ Check your understanding & A quick self-test before moving on. \\
			\textcolor{textgray}{\faBookOpen}\ Sources \& further reading & Curated references backing each module. \\
			\bottomrule
		\end{tabular}
		\caption{The callout boxes used throughout this guide.}
	\end{table}
	
	% =====================================================================
	\section*{Prerequisites and required background}
	\markright{Prerequisites and required background}
	\phantomsection
	\addcontentsline{toc}{section}{Prerequisites and required background}
	% =====================================================================
	
	The table below is the assumed starting point. None of the rows
	requires expert mastery; the goal of each is narrow and practical.
	
	\begin{table}[ht]
		\centering
		\small\sffamily
		\renewcommand{\arraystretch}{1.35}
		\rowcolors{2}{white}{primary!7}
		\begin{tabular}{@{}L{4.4cm} L{3.6cm} L{6.4cm}@{}}
			\toprule
			\rowcolor{primary}
			\textcolor{white}{\textbf{Skill}} &
			\textcolor{white}{\textbf{Required level}} &
			\textcolor{white}{\textbf{Primary use case}} \\
			\midrule
			Python programming & Intermediate & Implementing algorithms, data manipulation \\
			NumPy \& Pandas & Basic & Matrix operations and data cleaning \\
			Calculus \& Linear Algebra & Basic & Derivatives, partials, vectors, and matrices \\
			\bottomrule
		\end{tabular}
		\caption{Prerequisite matrix: the background assumed by this guide.}
		\label{tab:prereq}
	\end{table}
	
	\begin{table}[ht]
		\centering
		\small\sffamily
		\renewcommand{\arraystretch}{1.3}
		\caption{Notation reference table.}
		\begin{tabular}{@{}L{3.5cm} L{11.3cm}@{}}
			\toprule
			\rowcolor{primary}
			\textcolor{white}{\textbf{Symbol / Notation}} & \textcolor{white}{\textbf{Description and Typical Context}} \\
			\midrule
			$\mathbf{x}, \mathbf{y}, \mathbf{z}$ & Bold lowercase letters denote column vectors (e.g., sample features $\mathbf{x} \in \mathbb{R}^n$). \\
			$X, W, A$ & Uppercase letters denote matrices (e.g., dataset matrix $X \in \mathbb{R}^{m \times n}$, weights $W$). \\
			$\boldsymbol{\theta}, \mathbf{w}, \boldsymbol{\beta}$ & Model parameter vectors (e.g., regression weights $\boldsymbol{\beta}$, neural net parameters $\boldsymbol{\theta}$). \\
			$\mathcal{L}(\hat{y}, y),\; J(\boldsymbol{\theta})$ & Loss function for a single instance ($\mathcal{L}$) vs.\ cost function averaged over data ($J$). \\
			$\nabla_{\!\boldsymbol{\theta}} J$ & Gradient vector of partial derivatives with respect to parameters $\boldsymbol{\theta}$. \\
			$\Sigma$ vs.\ $\sum$ & Uppercase sigma $\Sigma$ denotes a Covariance Matrix; summation operator $\sum_{i=1}^m$ denotes summing a series. \\
			$\odot$ & Hadamard (element-wise) vector or matrix product. \\
			$\phi(z), \phi'(z)$ & Activation function and its scalar derivative. \\
			\bottomrule
		\end{tabular}
		\label{tab:notation}
	\end{table}
	
	\textbf{Python (intermediate).} You should be comfortable writing
	functions, loops, list comprehensions, and simple classes, and know your
	way around a virtual environment. The course is deliberately code-first:
	every mathematical idea is paired with a few lines of Python. A good
	refresher, free to read online, is \textcite{vanderplas2016}.
	
	\textbf{NumPy \& Pandas (basic).} You should be able to create arrays,
	slice them, and compute with them (\texttt{A @ B} for matrix products),
	and load a CSV into a DataFrame to clean and summarise it
	\parencite{harris2020, mckinney2010}. The practice boxes assume exactly
	this level and no more.
	
	\textbf{Calculus (basic).} You need the derivative as a slope, the power
	rule, and the chain rule. That is genuinely all --- gradients, the engine
	of all learning in this course, are just derivatives collected into a
	vector.
	
	\begin{keyconcept}[The two calculus ideas you actually need]
		\begin{itemize}
			\item \textbf{A derivative is a local slope:}
			$\;f'(x) \approx \dfrac{f(x+h)-f(x)}{h}$ --- it tells you which
			direction is \emph{downhill} and how steeply.
			\item \textbf{The chain rule composes slopes:}
			$\;\dfrac{d}{dx}\,f(g(x)) = f'(g(x))\,g'(x)$.
		\end{itemize}
		Backpropagation (Module~4) --- the efficient gradient-computation algorithm used to train modern
		neural networks --- is nothing more than this second rule applied
		systematically, layer by layer.
	\end{keyconcept}
	
	\textbf{Self-check before starting.}
	\begin{itemize}
		\item Can you write a Python function that returns the mean and standard
		deviation of a list of numbers?
		\item Given two NumPy arrays, can you compute their dot product and the
		shape of \texttt{A @ B}?
		\item Can you differentiate $f(w) = w^2$ and say what $f'(3)$ means
		geometrically?
	\end{itemize}
	If any answer is ``not yet'', spend a week with \textcite{vanderplas2016}
	and any standard calculus refresher; then return --- everything else is
	built for you from here.
	
	% =====================================================================
	\clearpage
	\section{Introduction and the AI Landscape}
	% =====================================================================
	
	\begin{modulecard}{\faGlobe\ Module 1 --- at a glance \hfill \small\normalfont (Estimated Study: 4 hours \textbullet{} Level: Beginner)}{primary}
		\textbf{Focus.} What artificial intelligence is (and is not), how the
		field evolved from hand-written rules to learning from data, and the
		infrastructure and ethical questions that shape modern practice.
		
		\medskip
		\textbf{Key concepts.} Defining intelligence; weak vs.\ strong AI; expert
		systems vs.\ data-driven AI; GPUs, TPUs, cloud ecosystems and data
		pipelines; bias, alignment, copyright and environmental impact.
		
		\medskip
		\textbf{After this module you will be able to} explain the difference
		between narrow and general AI, describe the paradigm shift from rules to
		data, sketch an end-to-end machine-learning pipeline, and name the main
		ethical questions attached to deploying an AI system.
	\end{modulecard}
	
	\subsection{Defining intelligence: weak vs.\ strong AI}
	
	Before we can build artificial intelligence we have to say what we mean
	by \emph{intelligence}. \textcite{russell2021} organise the many
	definitions along two axes: are we concerned with \emph{thought
		processes} or with \emph{behaviour}, and should the benchmark be
	\emph{human-like} performance or \emph{rationality} --- acting so as to
	achieve the best expected outcome? Alan Turing sidestepped the
	philosophy in 1950 with an operational test: if a machine's written
	answers are indistinguishable from a human's, is it not, for practical
	purposes, intelligent? \parencite{turing1950}. Modern AI research
	mostly aims at rational behaviour on well-defined tasks --- exactly what
	the learning algorithms of Modules~2--4 deliver.
	
	\begin{keyconcept}[Weak (narrow) AI vs.\ strong AI]
		\begin{center}
			\small\sffamily
			\renewcommand{\arraystretch}{1.25}
			\begin{tabular}{@{}L{6.6cm} L{7.9cm}@{}}
				\toprule
				\rowcolor{primary}
				\textcolor{white}{\textbf{Weak / narrow AI (today)}} &
				\textcolor{white}{\textbf{Strong AI / AGI (hypothetical)}} \\
				\midrule
				Excels at one well-defined task: spam filters, image classifiers, next-token prediction &
				Reasons, plans and learns across arbitrary domains, as a human can \\
				Learns from task-specific data; capability does not transfer &
				Would transfer knowledge flexibly to novel problems \\
				Every deployed system, including today's LLMs, is narrow in this sense &
				Unknown whether or when it is achievable; an active research and policy debate \\
				\bottomrule
			\end{tabular}
		\end{center}
	\end{keyconcept}
	
	A second classical debate is worth knowing. Philosopher John Searle's
	\emph{Chinese room} argument holds that symbol manipulation alone is
	not understanding: a person following rules could answer questions in
	Chinese without knowing any Chinese \parencite{searle1980}. The argument
	has returned with force in the era of large language models, because
	fluent output is cheap while grounded understanding is hard to verify
	--- a theme we revisit in the ethics spotlight below.
	
	\subsection{From expert systems to data-driven AI}
	
	The term \emph{artificial intelligence} was coined by John McCarthy in the 1955 proposal for the 1956 Dartmouth workshop. For roughly the next three decades the dominant paradigm was
	\emph{symbolic}: human knowledge was hand-encoded as if--then rules.
	Expert systems such as MYCIN encoded several hundred rules elicited
	from physicians and impressed clinicians in blind evaluations --- yet
	only inside their narrow domain. Rules are brittle, and extracting them
	from experts (the \emph{knowledge-acquisition bottleneck}) is slow and
	expensive. Deep Blue's 1997 victory over Garry Kasparov represented the high-water mark of this symbolic, search-based paradigm --- demonstrating immense raw evaluation power, yet remaining strictly bound by human-engineered evaluation heuristics.
	
	Machine learning inverts the pipeline: instead of programming rules, we
	collect examples of the desired behaviour and let an algorithm discover
	the rules. In \textcite{mitchell1997}'s classic formulation, a program
	\emph{learns} from experience $E$ with respect to task $T$ and
	performance measure $P$ if its performance at $T$, measured by $P$,
	improves with $E$. Deep learning pushes the inversion one step further
	by learning not only the rules but the \emph{features} they operate on
	\parencite{lecun2015}. Three results announced the new era: AlexNet cut
	the ImageNet top-5 error rate from roughly 26\% to 16\% by training a
	deep network on two consumer GPUs \parencite{krizhevsky2012}; AlphaGo
	defeated the world Go champion using deep networks combined with search
	\parencite{silver2016}; and transformer language models (Module~4)
	turned plain text prediction into a general-purpose technology.
	
	\begin{keyconcept}[The paradigm shift: from rules to data]
		\textbf{Classical programming:}\quad rules $+$ data $\longrightarrow$ answers\\
		\textbf{Machine learning:}\quad data $+$ answers $\longrightarrow$ rules (the model)
		
		\smallskip
		Instead of hand-crafting the logic, we specify an objective and let an
		algorithm search for parameters that meet it --- which is precisely why
		the mathematics of Modules~2 and~3 matters so much.
	\end{keyconcept}
	
	\begin{figure}[ht]
		\centering
		\begin{tikzpicture}[x=1cm, y=1cm, >=Stealth]
			\draw[primary, line width=2pt] (0,0) -- (14.3,0);
			\draw[primary, -{Stealth}, line width=2pt] (14.3,0) -- (14.75,0);
			\ms{0.55}{1950}{Turing's test}{0.18}{south}
			\ms{2.50}{1956}{Dartmouth: ``AI''}{-0.18}{north}
			\ms{4.45}{1980s}{expert systems}{0.18}{south}
			\ms{6.40}{1997}{Deep Blue}{-0.18}{north}
			\ms{8.35}{2012}{AlexNet}{0.18}{south}
			\ms{10.30}{2017}{Transformers}{-0.18}{north}
			\ms{12.25}{2020}{GPT-3}{0.18}{south}
			\ms{14.20}{2022}{ChatGPT}{-0.18}{north}
		\end{tikzpicture}
		\caption{A selective history: from the founding of the field to the generative-AI era.}
		\label{fig:timeline}
	\end{figure}
	
	\subsection{The modern stack: compute, cloud, and data pipelines}
	
	\textbf{Compute: GPUs and TPUs.} Training a modern model is,
	computationally, one thing: enormous numbers of matrix--vector
	products. A CPU optimises a handful of fast, sequential cores; a GPU
	packs thousands of simpler cores that perform these products in
	parallel. That is why \textcite{krizhevsky2012} could train a vision
	network that CPU-era attempts could not, and why hardware co-design now
	shapes chip roadmaps --- Google's Tensor Processing Units, for instance,
	execute the tensor algebra of neural networks directly
	\parencite{jouppi2017}. Memory bandwidth is the other half of the story:
	modern accelerators are as much about feeding data as about computing
	with it.
	
	\textbf{Cloud ecosystems.} The same hardware is available to anyone,
	rented by the hour: a laptop, a browser, and a credit card are enough
	to requisition a cluster of GPUs with a reproducible software image.
	This democratised experimentation --- and concentrated an enormous
	amount of AI capability in a handful of providers.
	
	\textbf{Data engineering, dataset construction, and contamination.} Models live inside a loop, not a vacuum
	(Figure~\ref{fig:pipeline}). Practitioners spend most project time collecting, cleaning, deduplicating, and validating data \parencite{mckinney2010}. Modern AI development requires rigorous data hygiene:
	\begin{itemize}[nosep]
		\item \textbf{Dataset Construction \& Labeling Noise:} Human annotation introduces errors, ambiguities, and labeler bias. Clean labels are often more valuable than model architecture tweaks.
		\item \textbf{Train/Test Leakage \& Contamination:} Pretraining datasets scraped at web scale frequently suffer from \emph{benchmark contamination} --- test set questions leaking into training corpora, inflating evaluation scores artificially. Temporal leakage (using future data to predict past events) and group leakage (splitting samples from the same patient/user across train and test sets) yield invalid accuracy claims.
		\item \textbf{Data Augmentation \& Synthetic Data:} Generating synthetic training examples or expanding datasets via domain-specific transforms mitigates data scarcity but risks amplifying generator biases.
	\end{itemize}
	
	\begin{practicebox}
\begin{lstlisting}
import pandas as pd
import numpy as np

# Ingest / synthetic creation for standalone demo
rng = np.random.default_rng(42)
df = pd.DataFrame({
    "created": pd.date_range("2026-01-01", periods=100, freq="h"),
    "amount": rng.uniform(10, 200, 100),
    "quantity": rng.integers(1, 5, 100),
    "region": rng.choice(["North", "South", "East", "West"], 100)
})
df = df.dropna(subset=["amount"])           # basic cleaning
df["revenue"] = df["amount"] * df["quantity"]
summary = df.groupby("region")["revenue"].agg(["count", "mean", "sum"])
print(summary)
\end{lstlisting}
		Pandas covers the ingest, clean, transform, aggregate loop
		that occupies most of a data scientist's day.
	\end{practicebox}
	
	\begin{figure}[ht]
		\centering
		\begin{tikzpicture}[
			font=\sffamily\scriptsize,
			stage/.style={draw=primary, fill=primary!8, rounded corners=1.8mm,
				minimum height=9mm, minimum width=2.55cm, align=center},
			>=Stealth]
			\node[stage] (a) {Data\\collection};
			\node[stage] (b) [right=0.45cm of a] {Cleaning \&\\preparation};
			\node[stage] (c) [right=0.45cm of b] {Model\\training};
			\node[stage] (d) [right=0.45cm of c] {Evaluation};
			\node[stage] (e) [right=0.45cm of d] {Deployment \&\\monitoring};
			\draw[->, thick, secondary] (a) -- (b);
			\draw[->, thick, secondary] (b) -- (c);
			\draw[->, thick, secondary] (c) -- (d);
			\draw[->, thick, secondary] (d) -- (e);
			\draw[->, thick, accent, dashed]
			(e.south) .. controls ($(e.south)+(0,-0.75)$) and ($(a.south)+(0,-0.75)$) .. (a.south);
			\node[font=\sffamily\scriptsize, text=accent, align=center]
			at ($(c.south)+(0,-0.78)$) {feedback: monitor, retrain on fresh data};
		\end{tikzpicture}
		\caption{A modern data pipeline: the machine-learning lifecycle as a loop, not a line.}
		\label{fig:pipeline}
	\end{figure}
	
	\subsection{Ethics, safety, and security: the non-optional module}
	
	Technical capability has outpaced governance. Six essential dimensions should
	accompany every system you build:
	
	\begin{ethicsbox}[Core dimensions of responsible AI and safety]
		\begin{description}[font=\sffamily\bfseries]
			\item[Bias \& Fairness.] Models inherit patterns from training data. Commercial classifiers have shown severe demographic accuracy gaps \parencite{buolamwini2018}. Crucially, \emph{simply stripping demographic features (e.g., race, gender) does NOT guarantee fairness}, because correlated proxy variables retain discriminatory signal. Furthermore, mathematical fairness criteria (such as Equalized Odds vs. Predictive Parity) frequently conflict under unequal base rates: \textbf{there is no single universally correct fairness metric}.
			\item[Privacy \& Security.] Training on sensitive records risks membership inference and model inversion attacks. Safeguards include Differential Privacy ($\epsilon, \delta$-guarantees) and federated aggregation. Additionally, deployment surfaces new threat vectors: \emph{prompt injection}, data poisoning, and model extraction.
			\item[Alignment \& Harm Mitigation.] Objective misalignment causes reward hacking \parencite{amodei2016, christian2020}. Red teaming, safety guardrails, and preference optimization (Module~4) help steer foundation models toward helpful and safe execution.
			\item[Copyright \& Provenance.] Data scraping raises unresolved legal questions surrounding fair use, licensing, watermarking, and artist attribution.
			\item[Environment.] Large-scale model training consumes significant energy and water \parencite{strubell2019, patterson2021}. Mitigations include right-sizing architectures and using energy-efficient accelerators.
			\item[Automation Bias \& Grounding.] \textcite{bender2021} warned that LLMs are \emph{stochastic parrots} reproducing surface-level statistics without grounded comprehension. Over-reliance (automation bias) leads users to accept unverified model claims.
		\end{description}
	\end{ethicsbox}
	
	\textbf{A practitioner's checklist.}
	\begin{itemize}
		\item Define the task, population served, and failure modes \emph{before} modeling.
		\item Audit training data for representation gaps, proxy correlations, and leakage.
		\item Evaluate metrics across demographic sub-groups; balance cost trade-offs.
		\item Document limitations via Model Cards, Datasheets, and clear system prompts.
		\item Establish red teaming, active monitoring, and human-in-the-loop oversight.
	\end{itemize}
	
	\begin{checkpoint}
		\begin{enumerate}
			\item A chess engine, a spam filter, and a claimed ``fully autonomous
			reasoning'' startup demo: classify each as weak or strong AI, and
			justify the label.
			\item Why did GPUs unlock deep learning at scale, whereas CPUs could not supply the necessary arithmetic throughput? Answer in terms of parallel hardware architecture.
			\item Your model screens job applications. Name two concrete harms and
			one mitigation for each.
		\end{enumerate}
	\end{checkpoint}
	
	\begin{sourcesbox}
		\small
		\begin{itemize}[nosep]
			\item \fullcite{russell2021}
			\item \fullcite{turing1950}
			\item \fullcite{searle1980}
			\item \fullcite{lecun2015}
			\item \fullcite{krizhevsky2012}
			\item \fullcite{silver2016}
			\item \fullcite{jouppi2017}
			\item \fullcite{buolamwini2018}
			\item \fullcite{amodei2016}
			\item \fullcite{christian2020}
			\item \fullcite{strubell2019}
			\item \fullcite{patterson2021}
			\item \fullcite{bender2021}
		\end{itemize}
	\end{sourcesbox}
	
	% =====================================================================
	\clearpage
	\section{Mathematics and Data Foundations}
	% =====================================================================
	
	\begin{modulecard}{\faRulerCombined\ Module 2 --- at a glance \hfill \small\normalfont (Estimated Study: 8 hours \textbullet{} Level: Intermediate)}{secondary}
		\textbf{Focus.} The small mathematical toolkit that quietly powers
		everything else in the course: linear algebra to \emph{represent} data,
		probability to \emph{quantify} uncertainty, and optimisation to
		\emph{learn} from mistakes.
		
		\medskip
		\textbf{Key concepts.} Vectors, matrices, eigenvalues and dimensionality
		reduction; Bayes' theorem, distributions and hypothesis testing; loss
		and cost functions, gradients, and gradient descent.
		
		\medskip
		\textbf{After this module you will be able to} describe any dataset as
		linear-algebra objects, reason correctly about conditional
		probabilities, and explain --- and implement --- gradient descent from
		scratch.
	\end{modulecard}
	
	Mathematics here is a working tool, not an exam subject. Each concept is
	introduced the moment it earns its keep, and each is paired with NumPy
	code that uses it.
	
	\subsection{Linear algebra: the language of data}
	
	Machine-learning data lives in vectors, matrices, and tensors. A house
	becomes a vector of measurements
	$\mathbf{x} = [\text{area}, \text{rooms}, \text{age}]$; a dataset of $m$ houses a matrix $X \in \mathbb{R}^{m \times n}$; a colour image a
	$H \times W \times 3$ tensor. The dot product measures alignment
	between vectors (its normalised version, cosine similarity, is the workhorse
	of recommendation systems and semantic search), and a matrix--vector
	product $W\mathbf{x}$ is a \emph{change of representation}: it rotates,
	stretches, and mixes coordinates. In Module~4 this becomes literal ---
	every neural-network layer \emph{is} one such matrix.
	
	\textbf{Eigenvalues and dimensionality reduction.} Eigenvectors are non-zero directions along which a linear transformation $A$ acts purely by scaling without changing direction:
	\begin{equation}
		A\mathbf{v} = \lambda \mathbf{v}.
	\end{equation}
	The covariance matrix $\Sigma$ of a dataset is symmetric, so its
	eigenvectors are mutually orthogonal, and the eigenvector with the
	largest eigenvalue points along the direction of greatest variance.
	Projecting the data onto the top eigenvectors is \emph{principal
		component analysis} (PCA):
	\begin{equation}
		\mathbf{v}_1 = \arg\max_{\|\mathbf{v}\| = 1}\ \mathbf{v}^{\!\top}\Sigma\,\mathbf{v},
	\end{equation}
	i.e.\ the first principal component is the top eigenvector of $\Sigma$.
	Dropping low-variance directions keeps most of the signal in far fewer
	coordinates --- compression, de-noising, and two-dimensional plots of
	hundred-dimensional data \parencite[Chapters 5--6]{strang2019}. NumPy makes all of it
	concrete, and its compiled vector operations are orders of magnitude
	faster than Python loops \parencite{harris2020}.
	
	\begin{practicebox}
\begin{lstlisting}
import numpy as np

X = np.random.default_rng(0).normal(size=(500, 20))
Xc = X - X.mean(axis=0)             # centre the columns
C  = (Xc.T @ Xc) / (X.shape[0]-1)   # covariance matrix
evals, evecs = np.linalg.eigh(C)    # symmetric solver
pcs = evecs[:, ::-1][:, :2]         # top-2 principal components
X2  = Xc @ pcs                      # 20-D data projected to 2-D
print(X2.shape, evals[::-1][:2])    # (500, 2) and top variances
\end{lstlisting}
	\end{practicebox}
	
	\subsection{Probability and statistics: the language of uncertainty}
	
	Real data is noisy and future observations are stochastic, so predictions in
	machine learning are fundamentally probabilistic statements.
	
	\textbf{Expectation, variance, and covariance.} For a random variable $X$, the \textbf{expected value} $\mathbb{E}[X] = \int x p(x)\,dx$ (or $\sum x_i P(x_i)$) measures central tendency. The \textbf{variance} quantifies dispersion:
	\begin{equation}
		\mathrm{Var}(X) = \mathbb{E}\bigl[(X - \mathbb{E}[X])^2\bigr] = \mathbb{E}[X^2] - (\mathbb{E}[X])^2.
	\end{equation}
	For two variables $X$ and $Y$, \textbf{covariance} measures joint linear variation $\mathrm{Cov}(X, Y) = \mathbb{E}\bigl[(X - \mathbb{E}[X])(Y - \mathbb{E}[Y])\bigr]$. When normalized by standard deviations $\sigma_X \sigma_Y$, it yields the Pearson correlation coefficient $\rho_{X,Y} \in [-1, 1]$.
	
	\textbf{Common distributions.} A small family covers foundational ML tasks: the \emph{Bernoulli} for binary outcomes $y \in \{0, 1\}$; the \emph{Categorical} for $K$-class labels and LLM tokens; and the \emph{Gaussian (Normal)} distribution for measurement noise:
	\begin{equation}
		p(x) = \frac{1}{\sigma\sqrt{2\pi}}\,
		\exp\!\left(-\frac{(x-\mu)^2}{2\sigma^2}\right).
	\end{equation}
	
	\textbf{Likelihood, MLE, and MAP estimation.} Parameter fitting estimates unknown distribution parameters $\boldsymbol{\theta}$ from data $\mathcal{D} = \{x_1, \dots, x_m\}$. The \textbf{Likelihood function} $L(\boldsymbol{\theta}) = p(\mathcal{D} \mid \boldsymbol{\theta}) = \prod_{i=1}^m p(x_i \mid \boldsymbol{\theta})$ evaluates data probability under candidate $\boldsymbol{\theta}$. Maximizing likelihood (or minimizing negative log-likelihood $-\ln L(\boldsymbol{\theta})$) gives the \textbf{Maximum Likelihood Estimator (MLE)}:
	\begin{equation}
		\hat{\boldsymbol{\theta}}_{\mathrm{MLE}} = \arg\max_{\boldsymbol{\theta}} \sum_{i=1}^{m} \ln p(x_i \mid \boldsymbol{\theta}).
	\end{equation}
	When prior knowledge $p(\boldsymbol{\theta})$ is available, \textbf{Maximum A Posteriori (MAP)} estimation incorporates the prior via Bayes' theorem:
	\begin{equation}
		\hat{\boldsymbol{\theta}}_{\mathrm{MAP}} = \arg\max_{\boldsymbol{\theta}} \left[ \sum_{i=1}^{m} \ln p(x_i \mid \boldsymbol{\theta}) + \ln p(\boldsymbol{\theta}) \right],
	\end{equation}
	which directly derives $\ell_2$ (Ridge) and $\ell_1$ (Lasso) regularisation terms from Gaussian and Laplace priors, respectively.
	
	\textbf{Bayes' theorem.} The foundation of probabilistic updating connects prior beliefs to observed data:
	\begin{equation}
		P(A \mid B) \;=\; \frac{P(B \mid A)\,P(A)}{P(B)}.
		\label{eq:bayes}
	\end{equation}
	
	\begin{keyconcept}[Bayes' theorem, in words]
		Posterior $\propto$ likelihood $\times$ prior: what you believe
		\emph{after} seeing evidence is proportional to how well each
		hypothesis would have produced that evidence, weighted by what you
		believed \emph{before}.
	\end{keyconcept}
	
	A medical screening example demonstrates why base rates dominate intuition. A diagnostic test with 99\% sensitivity and a 5\% false-positive rate, applied to a rare disease with 1\% prevalence, yields
	\begin{align}
		P(\text{disease}\mid +)
		&= \frac{P(+\mid\text{disease})\,P(\text{disease})}{P(+)}\\[2pt]
		&= \frac{\overbrace{0.99}^{\text{sensitivity}} \times \overbrace{0.01}^{P(\text{disease})}}{\underbrace{0.99}_{\text{sensitivity}} \times \underbrace{0.01}_{P(\text{disease})} \,+\, \underbrace{0.05}_{\text{false pos rate}} \times \underbrace{0.99}_{P(\text{no disease})}}
		\;\approx\; 0.17 :
	\end{align}
	only one positive in six is a true case, because the disease is rare to begin with. Note that the two $0.99$ terms in the denominator represent entirely different quantities: test sensitivity $P(+\mid\text{disease})=0.99$ versus healthy prior $P(\text{no disease})=1-0.01=0.99$. The same base-rate logic governs fraud alerts, security anomaly detection, and imbalanced classification.
	
	\textbf{Hypothesis testing, confidence intervals, and the CLT.} To discipline empirical findings:
	\begin{itemize}[nosep]
		\item \textbf{Central Limit Theorem (CLT):} For independent observations with finite variance $\sigma^2$, the standard error of the sample mean is exactly $\sigma/\sqrt{n}$. The CLT additionally states that as $n \to \infty$, the standardized sample mean distribution converges to a standard normal distribution $\mathcal{N}(0, 1)$.
		\item \textbf{Confidence Intervals \& Bootstrap Resampling:} Rather than relying solely on single-point estimates or $p$-values (which can be vulnerable to $p$-hacking under multiple testing), practitioners construct 95\% confidence intervals via asymptotic standard errors or non-parametric \emph{bootstrap resampling} (repeatedly sampling data with replacement to estimate metric variance empirical distributions).
	\end{itemize}
	
	\begin{practicebox}
\begin{lstlisting}
import numpy as np

rng = np.random.default_rng(42)
# 10,000 experiments: average of 40 Exponential(1) draws each
means = rng.exponential(scale=1.0, size=(10_000, 40)).mean(axis=1)
print(f"Sample Mean: {means.mean():.3f} | SE: {means.std():.3f}") # ~1.0 & 1/sqrt(40)
\end{lstlisting}
		The standard deviation of sample means matches standard error $\sigma/\sqrt{n}$, illustrating CLT convergence.
	\end{practicebox}
	
	\subsection{Optimisation: how models learn}
	
	A \textbf{loss function} $\mathcal{L}(\hat{y}, y)$ scores error on a single instance; the \textbf{cost function} $J(\boldsymbol{\theta})$ (empirical risk) averages loss over $m$ samples:
	\begin{equation}
		J(\boldsymbol{\theta}) = \frac{1}{2m}\sum_{i=1}^{m}
		\bigl(h_{\boldsymbol{\theta}}(\mathbf{x}_i) - y_i\bigr)^{2}.
	\end{equation}
	The \textbf{gradient} $\nabla_{\!\boldsymbol{\theta}} J$ points in the direction of steepest \emph{ascent}. Parameter updates step in the opposite direction:
	\begin{equation}
		\boldsymbol{\theta}_{t+1} \;=\;
		\boldsymbol{\theta}_t \;-\; \eta\,\nabla_{\!\boldsymbol{\theta}}
		J(\boldsymbol{\theta}_t),
		\label{eq:gd}
	\end{equation}
	where $\eta > 0$ is the learning rate. For strictly convex objectives (e.g., OLS linear regression with full-rank $X$ or $\ell_2$ regularisation), gradient descent converges to the unique global minimum. For non-convex neural network loss surfaces, stochastic mini-batch gradient descent (SGD) and adaptive optimizers reliably locate high-quality local minima.
	
	\textbf{Adam and AdamW optimisers.} Standard SGD uses a fixed step size. The \textbf{Adam} optimiser \parencite{kingma2015} maintains exponential moving averages of gradients $g_t$ (first moment $m_t$) and squared gradients $g_t^2$ (second moment $v_t$):
	\begin{equation}
		m_t = \beta_1 m_{t-1} + (1-\beta_1)g_t, \qquad
		v_t = \beta_2 v_{t-1} + (1-\beta_2)g_t^{2},
	\end{equation}
	\begin{equation}
		\hat{m}_t = \frac{m_t}{1-\beta_1^{t}}, \qquad
		\hat{v}_t = \frac{v_t}{1-\beta_2^{t}}, \qquad
		\boldsymbol{\theta}_{t+1} = \boldsymbol{\theta}_t - \frac{\eta\,\hat{m}_t}{\sqrt{\hat{v}_t} + \epsilon}.
	\end{equation}
	In modern deep learning and Transformer training, \textbf{AdamW} \parencite{loshchilov2019} is strongly preferred over standard Adam because it decouples $\ell_2$ weight decay regularization from adaptive gradient moment updates, preventing gradient scaling from inadvertently suppressing weight decay penalty terms.

	\begin{practicebox}[Adam optimiser step from scratch]
\begin{lstlisting}
import numpy as np

def adam_step(theta, grad, m, v, t, lr=0.01,
              beta1=0.9, beta2=0.999, eps=1e-8):
    m = beta1 * m + (1 - beta1) * grad
    v = beta2 * v + (1 - beta2) * (grad ** 2)
    m_hat = m / (1 - beta1 ** t)
    v_hat = v / (1 - beta2 ** t)
    theta_next = theta - (lr * m_hat) / (np.sqrt(v_hat) + eps)
    return theta_next, m, v
\end{lstlisting}
		Adam updates each parameter using bias-corrected first and second moment estimates.
	\end{practicebox}
	
	\begin{figure}[ht]
		\centering
		\begin{tikzpicture}[x=1.55cm, y=0.82cm, >=Stealth]
			\draw[textgray!60] (-2.75,0) -- (2.75,0)
			node[below, font=\scriptsize\sffamily, textgray] {$w$};
			\draw[textgray!60] (0,-0.15) -- (0,6.7)
			node[left, font=\scriptsize\sffamily, textgray] {$J(w)$};
			\draw[primary, very thick, domain=-2.55:2.55, smooth, samples=60]
			plot (\x, \x*\x);
			\foreach \a/\b in {2.2/1.54, 1.54/1.08, 1.08/0.755, 0.755/0.528,
				0.528/0.37, 0.37/0.259, 0.259/0.181}
			\draw[->, accent, thick] (\a, \a*\a) -- (\b, \b*\b);
			\foreach \w in {2.2, 1.54, 1.08, 0.755, 0.528, 0.37, 0.259, 0.181}
			\fill[secondary] (\w, \w*\w) circle (1.8pt);
			\node[font=\scriptsize\sffamily, anchor=west, text=primary] at (2.35, 4.45) {$w_0$};
			\node[font=\scriptsize\sffamily, anchor=west, text=primary] at (1.68, 2.40) {$w_1$};
			\node[font=\scriptsize\sffamily, anchor=west, text=primary] at (0.15, 5.55) {$J(w)=w^2$};
			\node[font=\scriptsize\sffamily, anchor=north, text=secondary] at (0, 0.10) {$w^{\star}=0$};
		\end{tikzpicture}
		\caption{Gradient descent on $J(w)=w^2$ with learning rate $\eta=0.15$:
			each step $w \leftarrow w - \eta\,2w$ shrinks the distance to the
			minimum by 30\%.}
		\label{fig:gd}
	\end{figure}
	
	\begin{practicebox}
\begin{lstlisting}
w, lr = 5.0, 0.15
for step in range(8):
    grad = 2 * w            # derivative of f(w) = w**2
    w = w - lr * grad       # the update rule, eq. (11)
    print(f"step {step}:  w = {w:.4f}")
    # w -> 0.2882: convergence toward the minimum
\end{lstlisting}
	\end{practicebox}
	
	\begin{checkpoint}
		\begin{enumerate}
			\item In plain words, what does an eigenvector of the covariance matrix
			tell you about a dataset?
			\item In the medical-test example, what happens to
			$P(\text{disease}\mid +)$ as prevalence rises to 10\%?
			\item What goes wrong, mechanically, when the learning rate in
			equation~\eqref{eq:gd} is set to $1.1$ for $J(w)=w^2$?
		\end{enumerate}
	\end{checkpoint}
	
	\begin{sourcesbox}
		\small
		\begin{itemize}[nosep]
			\item \fullcite{strang2019}
			\item \fullcite{wasserman2004}
			\item \fullcite{boyd2004}
			\item \fullcite{goodfellow2016}
			\item \fullcite{bishop2006}
			\item \fullcite{harris2020}
			\item \fullcite{mckinney2010}
			\item \fullcite{vanderplas2016}
		\end{itemize}
	\end{sourcesbox}
	
	% =====================================================================
	\clearpage
	\section{Classical Machine Learning: The Core Engine}
	% =====================================================================
	
	\begin{modulecard}{\faRobot\ Module 3 --- at a glance \hfill \small\normalfont (Estimated Study: 10 hours \textbullet{} Level: Intermediate)}{accent}
		\textbf{Focus.} The algorithms that still run the majority of
		real-world, structured-data workloads --- and the evaluation discipline
		that separates real learning from memorisation.
		
		\medskip
		\textbf{Key concepts.} Supervised learning: regression, decision trees,
		SVMs, ensembles. Unsupervised learning: K-Means clustering and anomaly
		detection. Reinforcement learning: MDPs, multi-armed bandits, exploration vs. exploitation. Evaluation: accuracy, precision, recall, F1, ROC-AUC.
		
		\medskip
		\textbf{After this module you will be able to} choose an appropriate
		classical algorithm for a task, cluster unlabelled data, solve multi-armed bandit problems, and --- most
		importantly --- evaluate any model honestly.
	\end{modulecard}
	
	On the tabular, structured data of business (customers, transactions,
	sensors), classical machine learning frequently still beats glamorous
	deep models. It is also the best place to internalise the habits ---
	splits, metrics, baselines --- that govern everything later
	\parencite{hastie2009, pedregosa2011}.
	
	\subsection{Supervised learning}
	
	Supervised learning starts from labelled pairs
	$\mathcal{D} = \{(\mathbf{x}_i, y_i)\}_{i=1}^{m}$ and a split into
	\emph{training}, \emph{validation}, and \emph{test} sets. If $y$ is a
	number the task is \emph{regression}; if $y$ is a category it is
	\emph{classification}.
	
	\textbf{Constructing the design matrix $X$ from raw data.} Consider a dataset of $m$ observations $\mathcal{D} = \{(\mathbf{x}_i, y_i)\}_{i=1}^{m}$, where each feature vector $\mathbf{x}_i = (x_{i1}, x_{i2}, \dots, x_{in})^{\!\top} \in \mathbb{R}^n$ represents $n$ input variables. A linear regression model expresses the prediction as an affine combination:
	\begin{equation}
		\hat{y}_i = \beta_0 + \beta_1 x_{i1} + \beta_2 x_{i2} + \dots + \beta_n x_{in} = \beta_0 + \boldsymbol{\beta}^{\!\top}\mathbf{x}_i,
	\end{equation}
	where $\beta_0 \in \mathbb{R}$ is the intercept (bias) and $\boldsymbol{\beta} = (\beta_1, \dots, \beta_n)^{\!\top} \in \mathbb{R}^n$ is the weight vector. To express predictions across all $m$ samples compactly as a single matrix--vector product $\hat{\mathbf{y}} = X\boldsymbol{\beta}_{\text{full}}$, we prepend a dummy constant feature $x_{i0} = 1$ to each sample vector $\tilde{\mathbf{x}}_i = (1, x_{i1}, \dots, x_{in})^{\!\top} \in \mathbb{R}^{n+1}$ and define the parameter vector $\boldsymbol{\beta}_{\text{full}} = (\beta_0, \beta_1, \dots, \beta_n)^{\!\top} \in \mathbb{R}^{n+1}$. Stacking these augmented row vectors $\tilde{\mathbf{x}}_i^{\!\top}$ forms the \textbf{design matrix} $X \in \mathbb{R}^{m \times (n+1)}$:
	\begin{equation}
		X = \begin{bmatrix}
			1 & x_{11} & x_{12} & \dots & x_{1n} \\
			1 & x_{21} & x_{22} & \dots & x_{2n} \\
			\vdots & \vdots & \vdots & \ddots & \vdots \\
			1 & x_{m1} & x_{m2} & \dots & x_{mn}
		\end{bmatrix} \in \mathbb{R}^{m \times (n+1)}.
	\end{equation}
	Under this representation, the Ordinary Least Squares (OLS) objective minimises the residual sum of squares $\|\mathbf{y} - X\boldsymbol{\beta}\|_2^2$, yielding the closed-form normal equations solution:
	\begin{equation}
		\hat{\boldsymbol{\beta}}_{\text{OLS}} = (X^{\!\top}X)^{-1}X^{\!\top}\mathbf{y},
	\end{equation}
	computed via QR factorisation or SVD for numerical stability \parencite{strang2019}.
	
	\textbf{Feature maps and basis expansion.} The matrix $X$ is not restricted to raw measurements. By applying a non-linear feature mapping $\boldsymbol{\phi}(\mathbf{x}) = [\phi_1(\mathbf{x}), \dots, \phi_d(\mathbf{x})]^{\!\top}$, linear models can capture complex non-linear curves. For example, a 1D input $x$ expanded into a degree-$d$ polynomial basis yields row vectors $\tilde{\mathbf{x}}_i^{\!\top} = [1, x_i, x_i^2, \dots, x_i^d]$, transforming non-linear fitting into a linear problem in feature space $\mathbb{R}^{d+1}$.
	
	\begin{practicebox}[Building Design Matrix $X$ \& Basis Expansion]
\begin{lstlisting}
import numpy as np

x_raw = np.array([1.5, 2.0, 3.5, 4.0, 5.5])
y = np.array([2.1, 2.9, 6.1, 8.2, 14.5])

# 1. Prepend bias column of 1s -> X in R^(m x 2)
X_linear = np.column_stack([np.ones(len(x_raw)), x_raw])

# 2. Basis expansion: Polynomial features [1, x, x^2, x^3]
X_poly = np.column_stack([x_raw**p for p in range(4)])

# Closed-form OLS via QR decomposition (R beta = Q^T y) & SVD (avoiding inv(X^T X)):
Q, R = np.linalg.qr(X_poly)
beta_poly_qr = np.linalg.solve(R, Q.T @ y)
beta_poly_svd = np.linalg.lstsq(X_poly, y, rcond=None)[0]

print(f"Linear X shape:   {X_linear.shape} | beta: {np.linalg.lstsq(X_linear, y, rcond=None)[0].round(2)}")
print(f"Poly X (QR):      {X_poly.shape} | beta: {beta_poly_qr.round(2)}")
print(f"Poly X (SVD):     {X_poly.shape} | beta: {beta_poly_svd.round(2)}")
\end{lstlisting}
	\end{practicebox}

	\textbf{Regularisation: Ridge ($\ell_2$) and Lasso ($\ell_1$).} When features are highly correlated or $n > m$, OLS exhibits high variance. Regularisation fights overfitting by penalising model complexity. Ridge regression adds an $\ell_2$ penalty:
	\begin{equation}
		\hat{\boldsymbol{\beta}}_{\text{Ridge}} = \arg\min_{\boldsymbol{\beta}}\; \frac{1}{2m}\|\mathbf{y} - X\boldsymbol{\beta}\|_2^{2} + \alpha \|\boldsymbol{\beta}\|_2^{2},
	\end{equation}
	which shrinks all coefficients smoothly toward zero but keeps them non-zero. By contrast, the \textbf{Lasso} ($\ell_1$ regularisation) penalises the sum of absolute values:
	\begin{equation}
		\hat{\boldsymbol{\beta}}_{\text{Lasso}} = \arg\min_{\boldsymbol{\beta}}\; \frac{1}{2m}\|\mathbf{y} - X\boldsymbol{\beta}\|_2^{2} + \alpha \|\boldsymbol{\beta}\|_1, \qquad \|\boldsymbol{\beta}\|_1 = \sum_{j=1}^{n} |\beta_j|.
	\end{equation}
	
	\textbf{Why Lasso induces sparsity.} Unlike $\ell_2$ sphere contours, the $\ell_1$ constraint set $\{\boldsymbol{\beta} : \|\boldsymbol{\beta}\|_1 \le t\}$ is a geometric polytope (a diamond in 2D) with sharp corners located directly on the coordinate axes. The elliptical level sets of the unconstrained OLS loss expand until they touch the constraint boundary; because the diamond's vertices lie on the axes where one or more coefficients equal zero, the solution frequently lands directly on a corner, driving uninformative feature weights strictly to zero \parencite{hastie2009}. Analytically, the absolute value function $|\beta_j|$ is non-differentiable at zero, yielding a subdifferential interval $[-1, 1]$. Coordinate descent solves for $\beta_j$ via the \emph{soft-thresholding operator}:
	\begin{equation}
		\mathcal{S}_{\gamma}(z) = \mathrm{sign}(z) \max\bigl(0,\; |z| - \gamma\bigr),
	\end{equation}
	which clamps small feature correlations below $\gamma = \alpha$ to exactly zero, performing automated feature selection.
	
	\begin{practicebox}[Lasso ($\ell_1$) Regularisation \& Numerical Solvers]
\begin{lstlisting}
import numpy as np

# Synthetic dataset with collinearity & sparse ground truth
rng = np.random.default_rng(42)
m, n = 100, 5
X = rng.normal(size=(m, n))
X[:, 4] = X[:, 0] + X[:, 1] + 1e-4 * rng.normal(size=m) # collinear feature
true_w = np.array([3.0, -2.5, 1.5, 0.0, 0.0])
y = X @ true_w + 0.2 * rng.normal(size=m)

# 1. Numerically stable OLS Normal Equations via QR & SVD
Q, R = np.linalg.qr(X)
beta_qr = np.linalg.solve(R, Q.T @ y)             # QR decomposition solver
beta_svd = np.linalg.lstsq(X, y, rcond=None)[0]   # SVD solver

# 2. Lasso via Coordinate Descent with Soft-Thresholding
def soft_threshold(z, gamma):
    return np.sign(z) * np.maximum(0, np.abs(z) - gamma)

def lasso_cd(X, y, alpha, max_iter=500):
    m, n = X.shape
    w = np.zeros(n)
    norm_x2 = np.sum(X**2, axis=0)
    for _ in range(max_iter):
        for j in range(n):
            r_j = y - (X @ w - X[:, j] * w[j])     # partial residual
            rho_j = X[:, j] @ r_j
            w[j] = soft_threshold(rho_j, m * alpha) / norm_x2[j]
    return w

w_lasso_np = lasso_cd(X, y, alpha=0.1)

print("True weights: ", true_w)
print("OLS QR:       ", beta_qr.round(2))          # unstable on collinear feature
print("NumPy Lasso:  ", w_lasso_np.round(2))       # zeroes noise & collinear weights
\end{lstlisting}
	\end{practicebox}

	\textbf{Classification: Sigmoid, Logistic Regression, and Softmax.} When target labels $y$ are categorical, linear output scores $z = \boldsymbol{\beta}^{\!\top}\mathbf{x} + \beta_0$ are unbounded and unsuitable for probabilities.
	
	\emph{Binary Logistic Regression.} For two classes $y \in \{0, 1\}$, the \textbf{sigmoid (logistic) function} maps any real score $z \in \mathbb{R}$ into a valid probability $p \in (0, 1)$:
	\begin{equation}
		\sigma(z) = \frac{1}{1 + e^{-z}} = \frac{e^{z}}{e^{z} + 1}.
	\end{equation}
	The model interprets $p = P(y=1 \mid \mathbf{x})$, implying that linear log-odds are given by $\ln\left(\frac{p}{1-p}\right) = \boldsymbol{\beta}^{\!\top}\mathbf{x} + \beta_0$. Fitting parameters via Maximum Likelihood Estimation yields the \textbf{Binary Cross-Entropy (BCE)} cost function:
	\begin{equation}
		\mathcal{L}(\boldsymbol{\beta}) = -\frac{1}{m} \sum_{i=1}^{m} \left[ y_i \ln \hat{y}_i + (1 - y_i) \ln (1 - \hat{y}_i) \right], \qquad \hat{y}_i = \sigma(\boldsymbol{\beta}^{\!\top}\mathbf{x}_i + \beta_0).
	\end{equation}
	
	\emph{Multinomial Logistic Regression and Softmax.} For $K > 2$ discrete classes $y \in \{1, \dots, K\}$, the model outputs a vector of unnormalised scores (logits) $\mathbf{z} = W\mathbf{x} + \mathbf{b} \in \mathbb{R}^K$, where $W \in \mathbb{R}^{K \times n}$. The \textbf{Softmax function} exponentiates and normalises logits into a categorical probability distribution $\hat{\mathbf{p}} = [\hat{p}_1, \dots, \hat{p}_K]^{\!\top}$:
	\begin{equation}
		\mathrm{softmax}(\mathbf{z})_k = \frac{e^{z_k}}{\sum_{j=1}^{K} e^{z_j}}, \qquad k = 1, \dots, K.
	\end{equation}
	Key properties of Softmax:
	\begin{enumerate}
		\item \textbf{Positivity and Normalisation:} Exponentiation enforces $e^{z_k} > 0$, and dividing by $\sum_j e^{z_j}$ guarantees $\sum_{k=1}^K \hat{p}_k = 1$ on the probability simplex.
		\item \textbf{Temperature Scaling:} Incorporating a temperature parameter $T > 0$ as $\mathrm{softmax}(\mathbf{z}/T)$ controls confidence: as $T \to 0^+$ predictions sharpen toward a one-hot argmax, whereas high $T \to \infty$ produces a uniform distribution (a mechanism central to LLM sampling in Module~4).
		\item \textbf{Generalisation of Sigmoid:} For $K=2$, Softmax on $[z_1, z_2]$ simplifies identically to Sigmoid on $(z_1 - z_2)$.
	\end{enumerate}
	Training multinomial logistic regression minimises the \textbf{Categorical Cross-Entropy} loss:
	\begin{equation}
		\mathcal{L}(W, \mathbf{b}) = -\frac{1}{m} \sum_{i=1}^{m} \sum_{k=1}^{K} y_{i,k} \ln \hat{p}_{i,k},
	\end{equation}
	where $y_{i,k} \in \{0, 1\}$ is the one-hot target indicator for class $k$.
	
	\begin{practicebox}[Sigmoid, Softmax \& Cross-Entropy Loss from Scratch]
\begin{lstlisting}
import numpy as np

def sigmoid(z):
    return 1 / (1 + np.exp(-np.clip(z, -500, 500)))

def softmax(Z, temp=1.0):
    # Subtract max logit per row for numerical stability
    Z_scaled = Z / temp
    exp_Z = np.exp(Z_scaled - np.max(Z_scaled, axis=-1, keepdims=True))
    return exp_Z / np.sum(exp_Z, axis=-1, keepdims=True)

# 3 samples, 4 classes (logits Z in R^(3 x 4))
logits = np.array([[2.0, 1.0, 0.1, 3.0],
                   [0.5, 4.0, 1.2, 0.0],
                   [-1.0, 0.0, 5.0, 0.2]])

probs = softmax(logits, temp=1.0)
y_true = np.array([3, 1, 2]) # True targets for each sample

# Categorical Cross-Entropy Loss: -log(p_true)
nll = -np.log(probs[np.arange(len(y_true)), y_true])
print("Probabilities:\n", probs.round(3))
print(f"Per-sample loss: {nll.round(3)} | Mean: {nll.mean():.3f}")
\end{lstlisting}
	\end{practicebox}
	
	\textbf{Naive Bayes and $k$-NN baselines.} Bridging directly from Bayes' theorem (equation~\ref{eq:bayes}), the \textbf{Naive Bayes classifier} computes the posterior probability of class $y \in \mathcal{Y}$ given features $\mathbf{x} = (x_1, \dots, x_n)$ by assuming features are conditionally independent given $y$:
	\begin{equation}
		\hat{y} = \arg\max_{y \in \mathcal{Y}}\; P(y \mid \mathbf{x}) = \arg\max_{y \in \mathcal{Y}}\; P(y) \prod_{j=1}^{n} P(x_j \mid y).
	\end{equation}
	Depending on data types, feature likelihoods $P(x_j \mid y)$ assume specific parametric distributions:
	\begin{itemize}[nosep]
		\item \textbf{Gaussian Naive Bayes} for continuous features: $P(x_j \mid y) = \frac{1}{\sqrt{2\pi\sigma_{y,j}^2}} \exp\left(-\frac{(x_j - \mu_{y,j})^2}{2\sigma_{y,j}^2}\right)$, where $\mu_{y,j}$ and $\sigma_{y,j}^2$ are class-conditional empirical means and variances.
		\item \textbf{Multinomial \& Bernoulli Naive Bayes} for discrete token counts or binary indicators in NLP.
	\end{itemize}
	To prevent zero-frequency probability collapse on unseen features during inference, \textbf{Laplace smoothing} adds a prior pseudo-count $\alpha \ge 1$: $\hat{P}(x_j \mid y) = \frac{N_{y, x_j} + \alpha}{N_y + \alpha \cdot d}$. In practice, products of small probabilities cause numerical underflow; computations sum log-likelihoods instead: $\ln P(y) + \sum_{j=1}^n \ln P(x_j \mid y)$.

	The \textbf{$k$-Nearest Neighbors ($k$-NN)} algorithm is a non-parametric instance-based baseline that stores all training observations $\mathcal{D}$. For a query point $\mathbf{x}$, it identifies the set $\mathcal{N}_k(\mathbf{x})$ of $k$ closest samples under a distance metric (e.g., Euclidean distance $d_2(\mathbf{x}, \mathbf{z}) = \|\mathbf{x} - \mathbf{z}\|_2$ or Manhattan distance $d_1(\mathbf{x}, \mathbf{z}) = \|\mathbf{x} - \mathbf{z}\|_1$) and assigns the majority class vote:
	\begin{equation}
		\hat{y} = \mathrm{mode}\bigl(\{y_i : \mathbf{x}_i \in \mathcal{N}_k(\mathbf{x})\}\bigr).
	\end{equation}
	Hyperparameter $k$ controls the bias--variance trade-off: $k=1$ creates complex, low-bias decision boundaries prone to overfitting noise, whereas large $k$ over-smooths class margins (high bias). While $k$-NN requires zero training compute, inference cost scales linearly as $O(m \cdot n)$ per query; spatial indexing trees (KD-Trees, Ball Trees) accelerate spatial queries in low-to-moderate dimensions ($n \ll 50$).

	\begin{practicebox}[Gaussian Naive Bayes from Scratch \& $k$-NN Baselines]
\begin{lstlisting}
import numpy as np
from sklearn.neighbors import KNeighborsClassifier

class GaussianNBScratch:
    def fit(self, X, y):
        self.classes = np.unique(y)
        self.priors = np.array([np.mean(y == c) for c in self.classes])
        self.means = np.array([X[y == c].mean(axis=0) for c in self.classes])
        self.vars = np.array([X[y == c].var(axis=0) + 1e-9 for c in self.classes])
        return self

    def predict(self, X):
        # Log likelihood for numerical stability: log P(y) + sum log P(x_j|y)
        log_probs = []
        for i, c in enumerate(self.classes):
            log_prior = np.log(self.priors[i])
            log_pdf = -0.5 * np.sum(np.log(2 * np.pi * self.vars[i]) + 
                                   ((X - self.means[i])**2) / self.vars[i], axis=1)
            log_probs.append(log_prior + log_pdf)
        return self.classes[np.argmax(np.array(log_probs), axis=0)]

# Synthetic binary dataset
rng = np.random.default_rng(42)
X = rng.normal(size=(200, 4))
y = (X[:, 0] + X[:, 1] > 0).astype(int)

gnb = GaussianNBScratch().fit(X[:150], y[:150])
knn = KNeighborsClassifier(n_neighbors=5).fit(X[:150], y[:150])

print(f"Gaussian NB (Scratch) Acc: {(gnb.predict(X[150:]) == y[150:]).mean():.3f}")
print(f"k-NN (k=5) Test Acc:       {knn.score(X[150:], y[150:]):.3f}")
\end{lstlisting}
		Naive Bayes computes class posterior probabilities in log-space via Gaussian feature likelihoods, matching $k$-NN baseline performance on smooth distributions.
	\end{practicebox}

	\textbf{Decision trees and ensembles.} A decision tree partitions feature space recursively into axis-aligned rectangular regions using a sequence of greedy binary splits. At node $m$, the algorithm searches over all features $j$ and split thresholds $t$ to maximize impurity reduction $\Delta I = I(m) - \frac{|S_L|}{|S|} I(S_L) - \frac{|S_R|}{|S|} I(S_R)$, where $|S|$ is node sample count. Common classification impurity metrics for class proportions $\mathbf{p} = (p_1, \dots, p_K)$ include:
	\begin{equation}
		\text{Gini Impurity:}\quad I_G(\mathbf{p}) = 1 - \sum_{k=1}^{K} p_k^2, \qquad \text{Entropy:}\quad H(\mathbf{p}) = -\sum_{k=1}^{K} p_k \log_2 p_k .
	\end{equation}
	For regression trees, splits minimize variance (MSE) $I_{\mathrm{MSE}}(S) = \frac{1}{|S|}\sum_{i \in S}(y_i - \bar{y}_S)^2$. Unconstrained trees grow until leaves are pure, causing severe overfitting (high variance). Overfitting is controlled via pre-pruning constraints (\texttt{max\_depth}, \texttt{min\_samples\_split}) or cost-complexity post-pruning $R_\alpha(T) = R(T) + \alpha |T|$.

	To overcome single-tree overfitting, \textbf{ensemble methods} combine multiple weak tree learners \parencite{hastie2009}:
	\begin{itemize}[nosep]
		\item \textbf{Random Forests (Bagging):} Grows $B$ independent de-correlated trees on bootstrap samples of training rows, evaluating only a random subset of $\sqrt{n}$ features per split \parencite{breiman2001}. Averaging tree outputs reduces variance by factor $\mathrm{Var}(\bar{X}) = \rho \sigma^2 + \frac{1-\rho}{B}\sigma^2$, where $\rho$ is pairwise tree correlation.
		\item \textbf{Gradient Boosting (XGBoost / LightGBM):} Builds trees sequentially: $\hat{y}_i^{(t)} = \hat{y}_i^{(t-1)} + \eta f_t(\mathbf{x}_i)$, where step $t$ fits tree $f_t$ to pseudo-residuals $r_{i,t} = -\left[\frac{\partial L(y_i, \hat{y})}{\partial \hat{y}}\right]_{\hat{y}^{(t-1)}}$ \parencite{chen2016}. Gradient boosting achieves state-of-the-art accuracy on tabular data.
	\end{itemize}

	\begin{practicebox}[Decision Trees, Gini Impurity \& Ensembles]
\begin{lstlisting}
import numpy as np
from sklearn.ensemble import RandomForestClassifier, HistGradientBoostingClassifier

def gini_impurity(y):
    """Compute Gini Impurity I_G(p) = 1 - sum(p_k^2)"""
    if len(y) == 0: return 0.0
    p = np.bincount(y) / len(y)
    return 1.0 - np.sum(p**2)

# Synthetic classification dataset with non-linear interaction
rng = np.random.default_rng(42)
X = rng.normal(size=(500, 10))
y = (X[:, 0]*2 - X[:, 1] + X[:, 2]**2 > 1.0).astype(int)

rf = RandomForestClassifier(n_estimators=100, max_depth=5, random_state=42).fit(X[:400], y[:400])
gb = HistGradientBoostingClassifier(max_iter=100, max_depth=5, random_state=42).fit(X[:400], y[:400])

print(f"Root Split Gini Impurity:   {gini_impurity(y[:400]):.3f}")
print(f"Random Forest Test Acc:     {rf.score(X[400:], y[400:]):.3f}")
print(f"Gradient Boosting Test Acc: {gb.score(X[400:], y[400:]):.3f}")
\end{lstlisting}
		Ensembles aggregate decision tree splits; gradient boosting iteratively corrects prior tree errors to model complex non-linear feature interactions.
	\end{practicebox}

	\textbf{Support Vector Machines (SVMs).} The SVM finds the decision boundary $f(\mathbf{x}) = \mathbf{w}^{\!\top}\mathbf{x} + b = 0$ maximizing the margin $\frac{2}{\|\mathbf{w}\|_2}$ to the nearest training points \parencite{cortes1995}. For non-linearly separable data, the \textbf{soft-margin formulation} introduces slack variables $\xi_i \ge 0$:
	\begin{equation}
		\min_{\mathbf{w}, b, \boldsymbol{\xi}}\; \frac{1}{2}\|\mathbf{w}\|_2^2 + C \sum_{i=1}^{m} \xi_i \qquad \text{s.t.}\quad y_i(\mathbf{w}^{\!\top}\mathbf{x}_i + b) \ge 1 - \xi_i, \quad \xi_i \ge 0,
	\end{equation}
	where target labels $y_i \in \{-1, +1\}$ and hyperparameter $C > 0$ trades off margin maximization against training errors. Mathematically, soft-margin SVM minimizes unconstrained \textbf{Hinge Loss} plus $\ell_2$ regularization: $\mathcal{L} = \sum_{i=1}^m \max\bigl(0, 1 - y_i f(\mathbf{x}_i)\bigr) + \frac{1}{2C}\|\mathbf{w}\|_2^2$.

	Formulating the dual problem using Lagrange multipliers $\alpha_i \ge 0$ yields:
	\begin{equation}
		\max_{\boldsymbol{\alpha}}\; \sum_{i=1}^{m} \alpha_i - \frac{1}{2}\sum_{i=1}^{m}\sum_{j=1}^{m} \alpha_i \alpha_j y_i y_j \bigl(\mathbf{x}_i^{\!\top}\mathbf{x}_j\bigr) \qquad \text{s.t.}\quad 0 \le \alpha_i \le C, \quad \sum_{i=1}^{m} \alpha_i y_i = 0.
	\end{equation}
	Observations with $\alpha_i > 0$ form the \textbf{support vectors} lying on or inside the margin. The dual formulation reveals that predictions depend exclusively on dot products $\mathbf{x}_i^{\!\top}\mathbf{x}_j$. Replacing this dot product with a non-linear kernel function $K(\mathbf{x}_i, \mathbf{x}_j) = \boldsymbol{\phi}(\mathbf{x}_i)^{\!\top}\boldsymbol{\phi}(\mathbf{x}_j)$ --- the \textbf{Kernel Trick} --- allows SVMs to construct non-linear decision boundaries in high-dimensional Hilbert spaces without explicitly computing expensive feature maps $\boldsymbol{\phi}(\cdot)$:
	\begin{itemize}[nosep]
		\item \textbf{Linear Kernel:} $K(\mathbf{x}, \mathbf{z}) = \mathbf{x}^{\!\top}\mathbf{z}$.
		\item \textbf{Polynomial Kernel:} $K(\mathbf{x}, \mathbf{z}) = (\mathbf{x}^{\!\top}\mathbf{z} + c)^d$.
		\item \textbf{Radial Basis Function (RBF / Gaussian Kernel):} $K(\mathbf{x}, \mathbf{z}) = \exp\bigl(-\gamma \|\mathbf{x} - \mathbf{z}\|_2^2\bigr)$, where $\gamma > 0$ controls the inverse length scale (higher $\gamma$ yields a smaller effective radius and tighter decision boundaries; lower $\gamma$ yields broader spatial influence).
	\end{itemize}

	\begin{practicebox}[Support Vector Machines: Linear vs.\ RBF Kernel]
\begin{lstlisting}
import numpy as np
from sklearn.svm import SVC

# Generate non-linearly separable concentric circles dataset
rng = np.random.default_rng(42)
X = rng.uniform(-2, 2, size=(300, 2))
y = ((X[:, 0]**2 + X[:, 1]**2) < 1.5).astype(int)

# Compare Linear vs RBF Kernel SVM
svm_lin = SVC(kernel='linear', C=1.0).fit(X[:240], y[:240])
svm_rbf = SVC(kernel='rbf', C=1.0, gamma='scale').fit(X[:240], y[:240])

print(f"Linear Kernel Test Acc: {svm_lin.score(X[240:], y[240:]):.3f}")
print(f"RBF Kernel Test Acc:    {svm_rbf.score(X[240:], y[240:]):.3f}")
print(f"RBF Support Vectors:    {len(svm_rbf.support_vectors_)} / 240 samples")
\end{lstlisting}
		The RBF kernel projects non-linearly separable points into implicit infinite-dimensional space, creating complex circular decision boundaries governed by support vectors.
	\end{practicebox}
	
	\begin{table}[ht]
		\centering
		\small\sffamily
		\renewcommand{\arraystretch}{1.3}
		\caption{Choosing a first algorithm for tabular data.}
		\begin{tabular}{@{}L{3.1cm} L{3.3cm} L{3.9cm} L{3.9cm}@{}}
			\toprule
			\rowcolor{accent}
			\textcolor{white}{\textbf{Algorithm}} &
			\textcolor{white}{\textbf{Reach for it when}} &
			\textcolor{white}{\textbf{Strengths}} &
			\textcolor{white}{\textbf{Watch out}} \\
			\midrule
			Linear / logistic regression & Baselines; small, clean data & Fast, interpretable, statistically principled & Underfits complex nonlinear structure \\
			Naive Bayes & Text classification, fast baseline & Very fast, handles high dimensions & Strong feature independence assumption \\
			$k$-NN & Small datasets, distance baselines & No training phase; flexible boundaries & Slow inference $O(mn)$; scale-sensitive \\
			Decision tree & Interpretability is the product & Human-readable rules; no feature scaling & Single trees overfit badly \\
			Random forest & Default robust choice & Variance reduction via bagging; few knobs & Less interpretable; larger memory \\
			Gradient boosting (XGBoost) & Maximum accuracy on tables & State of the art on structured data & Needs tuning; can overfit small data \\
			SVM & Small--medium data, many features & Max-margin generalisation; kernels & Scales poorly to very large datasets \\
			\bottomrule
		\end{tabular}
	\end{table}

	\textbf{Feature scaling requirement.} Algorithms that calculate spatial distances (K-Means, $k$-NN, SVM) or compute gradient steps (Logistic/Linear Regression, Neural Networks) are highly sensitive to feature scales: features with large magnitudes (e.g., house area in sq ft) dominate Euclidean metrics over features with small ranges (e.g., room count). Scaling is generally critical for distance-based algorithms ($k$-NN, SVM, K-Means) and often beneficial for gradient-based optimization (Logistic/Linear Regression, Neural Networks), while tree-based models are largely scale-invariant.
	
	\begin{ethicsbox}[CRITICAL: Avoid Data Leakage]
		\textbf{Fit scalers on training data ONLY.} A ubiquitous beginner error is fitting a scaler on the entire dataset prior to splitting. This causes \emph{data leakage}: information from the test set (mean $\mu$ and standard deviation $\sigma$) leaks into the training process, yielding artificially optimistic evaluation metrics. Always compute scaling parameters using \texttt{fit\_transform()} on the training split, and apply \texttt{transform()} on validation/test splits. Use \texttt{sklearn.pipeline.Pipeline} to automate this safely.
	\end{ethicsbox}

	\subsection{Unsupervised learning}
	
	Unsupervised learning discovers underlying structure without target labels.
	
	\textbf{K-Means \& Silhouette Score.} The canonical partitioning algorithm \textbf{K-Means} \parencite{macqueen1967} minimizes within-cluster inertia $J = \sum_{k=1}^K \sum_{\mathbf{x}_i \in C_k} \|\mathbf{x}_i - \boldsymbol{\mu}_k\|_2^2$ via Lloyd's algorithm (hard-assignment EM updates).\footnote{K-Means uses hard assignment (each point belongs to exactly one cluster), whereas Gaussian Mixture Models (GMMs) employ soft responsibilities via the full Expectation-Maximization algorithm as detailed below.} Cluster separation quality is quantified by the \textbf{Silhouette Score} $s(i) = \frac{b(i)-a(i)}{\max(a(i), b(i))}$ \parencite{rousseeuw1987}.
	
	\textbf{Density-based \& Mixture Models: DBSCAN \& GMMs.} K-Means assumes spherical, equal-sized clusters. To identify arbitrary non-convex clusters, \textbf{DBSCAN} \parencite{ester1996} groups points based on density ($\epsilon$-neighborhoods and minimum core samples), naturally isolating noise outliers. For probabilistic soft-clustering, \textbf{Gaussian Mixture Models (GMMs)} represent data distributions as a linear combination of $K$ Gaussians $\sum \pi_k \mathcal{N}(\mathbf{x} \mid \boldsymbol{\mu}_k, \Sigma_k)$, fitted via the Expectation-Maximization (EM) algorithm \parencite{bishop2006}.
	
	\textbf{Non-linear Dimensionality Reduction: t-SNE \& UMAP.} While PCA computes linear orthogonal projections, non-linear manifolds require non-linear embedding algorithms:
	\begin{itemize}[nosep]
		\item \textbf{t-SNE} \parencite{vanDerMaaten2008} converts pairwise Euclidean distances into conditional probabilities under Student-$t$ distributions, preserving local neighborhood structures in 2D/3D visualizations.
		\item \textbf{UMAP} \parencite{mcinnes2018} constructs fuzzy simplicial sets to optimize low-dimensional layouts, preserving both local and global structure with superior scalability.
	\end{itemize}

	\textbf{Anomaly Detection.} Methods include Isolation Forest \parencite{liu2008} (random feature splits isolating anomalies near root nodes), Elliptic Envelope, and Autoencoders (Module 4).
	
	\begin{practicebox}[DBSCAN, K-Means \& Silhouette Score]
\begin{lstlisting}
from sklearn.cluster import KMeans, DBSCAN
from sklearn.metrics import silhouette_score
from sklearn.datasets import make_blobs
import numpy as np

# Synthetic data generation for standalone execution
X_tr, _ = make_blobs(n_samples=150, centers=3, cluster_std=0.6, random_state=42)

# Fit K-Means and DBSCAN on synthetic data
kmeans = KMeans(n_clusters=3, n_init=10, random_state=42).fit(X_tr)
dbscan = DBSCAN(eps=0.5, min_samples=5).fit(X_tr)

score = silhouette_score(X_tr, kmeans.labels_)
print(f"K-Means Silhouette Score: {score:.3f}")
print(f"DBSCAN Clusters Found:   {len(set(dbscan.labels_)) - (1 if -1 in dbscan.labels_ else 0)}")
\end{lstlisting}
	\end{practicebox}

	\subsection{Reinforcement learning}
	
	While supervised learning requires target pairs $(\mathbf{x}_i, y_i)$ and unsupervised learning discovers unlabelled structure, \textbf{Reinforcement Learning (RL)} optimizes sequential decision-making through trial-and-error environment interaction \parencite{sutton2018}. An agent observes state $s_t \in \mathcal{S}$, executes action $a_t \in \mathcal{A}$ under policy $\pi(a \mid s)$, transitions to $s_{t+1}$, and receives reward $r_{t+1}$.
	
	\textbf{MDPs \& Bellman Equations.} Environments obey the Markov property: $P(S_{t+1}=s', R_{t+1}=r \mid S_t, A_t, \dots) = P(S_{t+1}=s', R_{t+1}=r \mid S_t, A_t)$. An MDP $(\mathcal{S}, \mathcal{A}, \mathcal{P}, \mathcal{R}, \gamma)$ defines state/action spaces, transition dynamics $\mathcal{P}(s' \mid s, a)$, reward function $\mathcal{R}$, and discount factor $\gamma \in [0, 1)$. Value functions $V^\pi(s) = \mathbb{E}_\pi[G_t \mid S_t=s]$ and $Q^\pi(s, a) = \mathbb{E}_\pi[G_t \mid S_t=s, A_t=a]$ satisfy the \textbf{Bellman Optimality Equation}:
	\begin{equation}
		Q^{\star}(s, a) = \sum_{s' \in \mathcal{S}} \mathcal{P}(s' \mid s, a) \left[ \mathcal{R}(s, a, s') + \gamma \max_{a' \in \mathcal{A}} Q^{\star}(s', a') \right].
	\end{equation}

	\textbf{Temporal Difference (TD) Learning, SARSA, \& Q-Learning.} Model-free RL learns directly from interaction:
	\begin{itemize}[nosep]
		\item \textbf{Monte Carlo vs. TD Learning:} Monte Carlo updates values after full episode completion, whereas TD learning updates predictions incrementally from immediate reward plus estimated next-state value (bootstrapping): $V(s_t) \leftarrow V(s_t) + \alpha [r_{t+1} + \gamma V(s_{t+1}) - V(s_t)]$.
		\item \textbf{SARSA (On-Policy) vs. Q-Learning (Off-Policy):} SARSA updates Q-values using the actual next action executed by the policy: $Q(s_t,a_t) \leftarrow Q(s_t,a_t) + \alpha [r_{t+1} + \gamma Q(s_{t+1}, a_{t+1}) - Q(s_t,a_t)]$. Off-policy Q-learning \parencite{watkins1992} updates using the maximum possible next-state action value $\max_a Q(s_{t+1}, a)$.
	\end{itemize}

	\textbf{Actor-Critic Methods \& Advantage Functions.} To reduce policy gradient variance, \textbf{Actor-Critic} architectures maintain an \emph{Actor} (policy $\pi_{\boldsymbol{\theta}}(a \mid s)$) and a \emph{Critic} (value function $V_{\boldsymbol{\phi}}(s)$). The update scales gradients by the \textbf{Advantage function}:
	\begin{equation}
		A(s, a) = Q(s, a) - V(s) \;\approx\; r_{t+1} + \gamma V(s_{t+1}) - V(s_t),
	\end{equation}
	quantifying how much better action $a$ is compared to the average expected value of state $s$. Modern algorithms like Proximal Policy Optimization (PPO) \parencite{schulman2017} clip policy updates to prevent destabilizing gradient steps, serving as the core engine for RLHF (Module 4).
	
	\textbf{Clarification on AlphaGo \& MCTS.} Note that Monte Carlo Tree Search (MCTS) in AlphaGo \parencite{silver2016} performs lookahead search over a perfect game simulator, combining heuristic rollouts with neural evaluation. This lookahead search over a simulator is distinct from model-based RL algorithms that must explicitly learn an empirical transition probability model $\hat{\mathcal{P}}(s' \mid s, a)$ from raw environment data.

	\subsection{Modern evaluation methodology \& diagnostics}
	
	Evaluating a model requires moving beyond simple train/test accuracy.
	
	\textbf{Classification Metrics: Confusion Matrix, PR-AUC, and ROC-AUC.} Precision $\frac{TP}{TP+FP}$, Recall $\frac{TP}{TP+FN}$, $F_1 = 2 \frac{P \cdot R}{P + R}$, and ROC-AUC measure classification performance. For severely imbalanced datasets (e.g., fraud, rare diseases), ROC-AUC can present an overly optimistic score because false positives are divided by a vast number of true negatives ($TN$). The \textbf{Precision-Recall Curve (PR-AUC)} is strongly preferred for imbalanced evaluation because it focuses exclusively on the positive class.
	
	\textbf{Probabilistic Calibration \& Proper Scoring Rules.} Model probabilities $\hat{p} \in [0, 1]$ should match real-world frequencies: among events assigned $\hat{p} = 0.8$, exactly 80\% should occur. \textbf{Calibration curves} (reliability diagrams) plot predicted probability bins against empirical frequencies. Quantitative calibration error is measured by the \textbf{Brier Score} $\frac{1}{m}\sum (\hat{y}_i - y_i)^2$ and Categorical Cross-Entropy (Log-Loss), both of which are \emph{proper scoring rules} minimized strictly when predicted probabilities equal true posterior distributions.
	
	\textbf{Advanced Resampling \& Data Drift.} Standard $k$-fold cross-validation fails on structured data dependencies:
	\begin{itemize}[nosep]
		\item \textbf{Nested Cross-Validation:} Outer folds evaluate generalization error while inner folds perform hyperparameter optimization, preventing optimistic hyperparameter selection bias.
		\item \textbf{Time-Series \& Grouped CV:} Time-dependent data requires rolling-window splits ($t_{\text{train}} < t_{\text{test}}$) to prevent temporal leakage; grouped data (multiple samples per user) requires \texttt{GroupKFold} to prevent user leakage across splits.
		\item \textbf{Data Drift \& Concept Drift:} In production, input feature distributions change over time (\emph{covariate shift} $P(X)$ changes) or target relationships alter (\emph{concept drift} $P(Y \mid X)$ changes), requiring monitoring statistics (Kolmogorov-Smirnov tests, Population Stability Index) and scheduled model retraining.
	\end{itemize}
	
	\begin{checkpoint}
		\begin{enumerate}
			\item For a cancer-screening tool, which metric would you prioritise,
			and what is the cost of that choice?
			\item Why does K-Means struggle with two concentric rings of points,
			and what does that reveal about its assumptions?
			\item In a $K$-armed bandit problem, how does increasing $\epsilon$ in an $\epsilon$-greedy strategy impact short-term vs.\ long-term return?
			\item Your cross-validation scores vary from 0.62 to 0.91 across folds.
			What does that variance tell you, and what would you do next?
		\end{enumerate}
	\end{checkpoint}
	
	\begin{sourcesbox}
		\small
		\begin{itemize}[nosep]
			\item \fullcite{mitchell1997}
			\item \fullcite{hastie2009}
			\item \fullcite{bishop2006}
			\item \fullcite{cortes1995}
			\item \fullcite{breiman2001}
			\item \fullcite{chen2016}
			\item \fullcite{macqueen1967}
			\item \fullcite{liu2008}
			\item \fullcite{sutton2018}
			\item \fullcite{watkins1992}
			\item \fullcite{mnih2015}
			\item \fullcite{schulman2017}
			\item \fullcite{silver2016}
			\item \fullcite{amodei2016}
			\item \fullcite{fawcett2006}
			\item \fullcite{pedregosa2011}
		\end{itemize}
	\end{sourcesbox}
	
	% =====================================================================
	\clearpage
	\section{Transition to Deep Learning and Modern AI}
	% =====================================================================
	
	\begin{modulecard}{\faBrain\ Module 4 --- at a glance \hfill \small\normalfont (Estimated Study: 8 hours \textbullet{} Level: Intermediate--Advanced)}{dpurple}
		\textbf{Focus.} The leap from engineered features to learned representation hierarchies: perceptrons, activation functions, backpropagation, regularization \& normalization, autoencoders, latent spaces, CNNs, RNNs/LSTMs, and self-attention in Transformers.
		
		\medskip
		\textbf{Key concepts.} Perceptrons, backpropagation, activation functions, vanishing gradients, Dropout, Batch/Layer Normalization, autoencoders, latent spaces, Word2Vec, CNNs, translation equivariance, RNNs, LSTMs, Scaled Dot-Product Attention, and Multi-Head Attention.
		
		\medskip
		\textbf{After this module you will be able to} explain how neural networks learn continuous latent representations, derive backpropagation matrix calculus, contrast normalization techniques, and build Multi-Head Attention mechanisms from scratch.
	\end{modulecard}
	
	The jump from classical ML to deep learning is a jump in \emph{who
		designs the features}. A classical pipeline is given hand-crafted
	inputs; a deep network learns a hierarchy of them --- edges, then
	textures, then parts, then objects \parencite{lecun2015}. What
	engineering achieved by hand, depth achieves with data and compute.
	
	\subsection{Neural networks: perceptrons, activations, backpropagation}
	
	The \textbf{perceptron} \parencite{rosenblatt1958} is a single neuron:
	a weighted sum passed through a threshold,
	\begin{equation}
		\hat{y} = \mathbb{1}\{\mathbf{w}^{\!\top}\mathbf{x} + b > 0\} \in \{0, 1\}
		\quad (\text{or }\mathrm{sgn}(\mathbf{w}^{\!\top}\mathbf{x} + b) \in \{-1, +1\}).
	\end{equation}
	It can learn any function that a straight line can separate --- and
	provably cannot learn XOR, a limitation that chilled the field's first
	winter \parencite{minsky1969}. The repair is to \emph{stack} units and
	interleave \textbf{nonlinear activation functions} between layers:
	\begin{equation}
		\mathbf{a}^{(\ell)} = \phi\!\left(W^{(\ell)}\,\mathbf{a}^{(\ell-1)} +
		\mathbf{b}^{(\ell)}\right),
		\qquad \ell = 1, \dots, L .
	\end{equation}
	Without $\phi$, the stack would collapse into a single linear map. The
	rectified linear unit, $\phi(z) = \max(0, z)$, became the default for
	being cheap, non-saturating, and sparse in its activations
	\parencite{nair2010}. Such networks are universal approximators
	\parencite{hornik1989} --- a statement about \emph{capacity}, with no
	guarantee that training finds good parameters.
	
	\begin{figure}[ht]
		\centering
		\begin{tikzpicture}[
			x=1.35cm, y=0.9cm,
			inp/.style={circle,  draw=primary,   fill=primary!12,   minimum size=6mm, inner sep=0pt},
			hid/.style={circle,  draw=dpurple,   fill=dpurple!12,   minimum size=6mm, inner sep=0pt},
			outp/.style={circle, draw=secondary, fill=secondary!12, minimum size=6mm, inner sep=0pt}]
			\foreach \i in {1,2,3}   \node[inp]  (I\i)  at (0,    \i-2)   {};
			\foreach \i in {1,...,4} \node[hid]  (H1\i) at (1.45, \i-2.5) {};
			\foreach \i in {1,...,4} \node[hid]  (H2\i) at (2.9,  \i-2.5) {};
			\foreach \i in {1,2}     \node[outp] (O\i)  at (4.35, \i-1.5) {};
			\foreach \i in {1,2,3} \foreach \j in {1,...,4}
			\draw[primary!35]   (I\i)  -- (H1\j);
			\foreach \i in {1,...,4} \foreach \j in {1,...,4}
			\draw[dpurple!35]   (H1\i) -- (H2\j);
			\foreach \i in {1,...,4} \foreach \j in {1,2}
			\draw[secondary!40] (H2\i) -- (O\j);
			\node[font=\sffamily\scriptsize\bfseries, text=primary,   align=center] at (0,    1.8) {input};
			\node[font=\sffamily\scriptsize\bfseries, text=dpurple,   align=center] at (1.45, 2.3) {hidden};
			\node[font=\sffamily\scriptsize\bfseries, text=dpurple,   align=center] at (2.9,  2.3) {hidden};
			\node[font=\sffamily\scriptsize\bfseries, text=secondary, align=center] at (4.35, 1.3) {output};
		\end{tikzpicture}
		\caption{A fully connected $3$--$4$--$4$--$2$ network. Every edge
			carries a learnable weight; each layer applies its matrix, then a
			nonlinearity.}
		\label{fig:nn}
	\end{figure}
		
	\begin{figure}[ht]
		\centering
		\begin{tikzpicture}
			\begin{axis}[
				width=0.72\textwidth, height=5.6cm,
				axis lines=middle,
				xlabel={$z$}, ylabel={$\phi(z)$},
				xmin=-3.3, xmax=3.3, ymin=-1.4, ymax=2.4,
				domain=-3.2:3.2, samples=101,
				tick label style={font=\scriptsize},
				label style={font=\small},
				legend style={at={(0.02,0.98)}, anchor=north west,
					legend cell align=left, font=\scriptsize\sffamily,
					draw=none, fill=none}]
				\addplot[primary, very thick]            {1/(1+exp(-x))};
				\addlegendentry{sigmoid}
				\addplot[secondary, very thick, dashed]  {2/(1+exp(-2*x))-1};
				\addlegendentry{tanh}
				\addplot[accent, very thick, domain=0:2.35] {max(0,x)};
				\addlegendentry{ReLU}
			\end{axis}
		\end{tikzpicture}
		\caption{Common activation functions. ReLU's lack of saturation for
			$z>0$ helped make deep networks trainable at scale \parencite{nair2010}.}
		\label{fig:activations}
	\end{figure}
	
	\textbf{Backpropagation} \parencite{rumelhart1986} computes the
	gradient of the loss $\mathcal{L}$ with respect to every weight 
	(the chain rule from Module~1) organised backwards through the layers. With
	$\delta^{(\ell)} \equiv \partial \mathcal{L}/\partial
	\mathbf{z}^{(\ell)}$:
	\begin{align}
		\boldsymbol{\delta}^{(L)}
		&= \nabla_{\mathbf{a}^{(L)}}\,\mathcal{L}\ \odot\
		\phi'\!\bigl(\mathbf{z}^{(L)}\bigr),\\
		\boldsymbol{\delta}^{(\ell)}
		&= \bigl(W^{(\ell+1)\top}\boldsymbol{\delta}^{(\ell+1)}\bigr)\ \odot\
		\phi'\!\bigl(\mathbf{z}^{(\ell)}\bigr),
		\qquad \ell = L-1, \dots, 1,\\
		\frac{\partial \mathcal{L}}{\partial W^{(\ell)}}
		&= \boldsymbol{\delta}^{(\ell)}\,
		\mathbf{a}^{(\ell-1)\top},
	\end{align}
	where $\odot$ is element-wise multiplication. Backprop is a
	gradient-\emph{computation} algorithm; the learning itself is still
	gradient descent (equation~\eqref{eq:gd}), now run on mini-batches
	through frameworks such as PyTorch or JAX.

	\textbf{Vanishing and exploding gradients.} In deep multi-layer architectures, backpropagating loss gradients across $L$ layers involves repeated matrix products backwards:
	\begin{equation}
		\boldsymbol{\delta}^{(1)} = \left( \prod_{\ell=2}^{L} W^{(\ell)\top} \mathrm{diag}\!\left(\phi'\!\left(\mathbf{z}^{(\ell)}\right)\right) \right) \boldsymbol{\delta}^{(L)}.
	\end{equation}
	Backpropagated signals scale with matrix singular values and activation derivatives $\phi'$. In naive unregularized deep stacks, if effective singular values or activation slopes remain consistently below 1, gradients decay exponentially to zero ($\boldsymbol{\delta}^{(1)} \to \mathbf{0}$, \emph{vanishing gradients}), leaving early layers untrained. Conversely, if effective products exceed 1, gradients explode exponentially ($\boldsymbol{\delta}^{(1)} \to \infty$, \emph{exploding gradients}). In modern deep networks, gradient dynamics depend strongly on architecture, initialization, normalization, and activation distributions. Deep networks employ explicit safeguards --- residual skip connections \parencite{he2016}, gradient clipping, Layer Normalization, Xavier/He weight initialization \parencite{glorot2010, he2015}, and non-saturating activations (ReLU) \parencite{nair2010} --- to maintain stable gradient norm propagation.

	\textbf{Regularization and normalization in deep networks.} Beyond basic weight decay, deep network training relies on specialized techniques to prevent overfitting and stabilize activation distributions across layers:
	
	\begin{table}[ht]
		\centering
		\small\sffamily
		\renewcommand{\arraystretch}{1.3}
		\rowcolors{2}{white}{primary!7}
		\begin{tabular}{@{}L{3.2cm} L{4.5cm} L{6.8cm}@{}}
			\toprule
			\rowcolor{primary}
			\textcolor{white}{\textbf{Technique}} & \textcolor{white}{\textbf{Mechanism \& Scope}} & \textcolor{white}{\textbf{Primary Benefit \& Use Case}} \\
			\midrule
			\textbf{Dropout} & Randomly zeroes activations with probability $p$ during training & Prevents co-adaptation; acts as an ensemble over sub-networks \\
			\textbf{Batch Norm} & Normalizes activations across the mini-batch dimension & Reduces internal covariate shift in CNNs; sensitive to batch size \\
			\textbf{Layer Norm} & Normalizes activations across feature channels per sample & Essential for Transformers and RNNs; independent of batch size \\
			\textbf{Weight Decay vs $\ell_2$} & Decouples coefficient penalty from adaptive moments (AdamW) & Prevents gradient scaling from eroding regularisation in LLMs \\
			\bottomrule
		\end{tabular}
		\caption{Comparison of regularization and normalization mechanisms in modern deep architectures.}
		\label{tab:reg_norm}
	\end{table}

	\begin{practicebox}[Two-Layer MLP with Manual Backpropagation on Non-Linear Data]
\begin{lstlisting}
import numpy as np

# Non-linear XOR / concentric circle dataset
rng = np.random.default_rng(42)
X_b = rng.uniform(-2, 2, size=(200, 2))
y_b = ((X_b[:, 0]**2 + X_b[:, 1]**2) < 1.2).astype(int).reshape(-1, 1)

# Initialize weights with He Normalization (scaling by sqrt(2/fan_in))
W1 = rng.normal(0, np.sqrt(2/2), size=(2, 8)); b1 = np.zeros((1, 8))
W2 = rng.normal(0, np.sqrt(2/8), size=(8, 1)); b2 = np.zeros((1, 1))

for epoch in range(500):
    # Forward pass: Input -> Hidden (ReLU) -> Output (Sigmoid)
    z1 = X_b @ W1 + b1; a1 = np.maximum(0, z1)   # ReLU activation
    z2 = a1 @ W2 + b2;  hat_y = 1 / (1 + np.exp(-np.clip(z2, -500, 500)))
    
    # Backward pass (Gradient computation via chain rule)
    dz2 = (hat_y - y_b) / len(X_b)               # dL/dz2 for BCE loss
    dW2 = a1.T @ dz2; db2 = dz2.sum(axis=0, keepdims=True)
    da1 = dz2 @ W2.T;  dz1 = da1 * (z1 > 0)      # ReLU derivative
    dW1 = X_b.T @ dz1; db1 = dz1.sum(axis=0, keepdims=True)
    
    # Gradient descent parameter updates
    W1 -= 0.1 * dW1; b1 -= 0.1 * db1
    W2 -= 0.1 * dW2; b2 -= 0.1 * db2

bce_loss = -np.mean(y_b*np.log(hat_y+1e-7) + (1-y_b)*np.log(1-hat_y+1e-7))
acc = np.mean((hat_y > 0.5) == y_b)
print(f"Epoch 500 | BCE Loss: {bce_loss:.4f} | Accuracy: {acc:.3f}")
\end{lstlisting}
	\end{practicebox}
	
	\subsection{Representation learning: from manual engineering to latent spaces}
	
	In classical machine learning (Module~3), algorithm accuracy hinges on manual feature engineering --- crafting input matrix columns $\mathbf{x} \in \mathbb{R}^n$ using domain heuristics (e.g., SIFT descriptors in vision, TF-IDF in NLP, polynomial features). Manual feature engineering is laborious, task-specific, fragile, and fails to scale to high-dimensional unstructured inputs like pixels or continuous audio waveforms. \textbf{Representation learning} \parencite{bengio2013} eliminates manual feature engineering: neural networks discover parameterized continuous transformations $\mathbf{z} = f_{\boldsymbol{\theta}}(\mathbf{x}) \in \mathbb{R}^d$ that automatically extract latent semantic structures and data-generating invariants.
	
	\textbf{The Manifold Hypothesis: linear vs.\ non-linear projections.} High-dimensional real-world data vectors $\mathbf{x} \in \mathbb{R}^n$ (such as $256 \times 256 \times 3$ RGB images in $\mathbb{R}^{196,608}$) occupy an infinitesimal fraction of ambient space. The \textbf{Manifold Hypothesis} posits that real-world data concentrates near a smooth, low-dimensional non-linear manifold $\mathcal{M} \subset \mathbb{R}^n$ of intrinsic dimension $d \ll n$. Linear dimensionality reduction (PCA, Module~2) computes orthogonal projections $\mathbf{z} = W^{\!\top}\mathbf{x}$, which capture global variance but fail to unfold non-linear manifolds (e.g., Swiss Roll geometry). Interleaving non-linear activation functions across hidden layers allows deep networks to learn non-linear coordinate charts, flattening complex curved manifolds in input space into linearly separable representations in latent space $\mathcal{Z}$.
	
	\begin{keyconcept}[The Core Principle of Representation Learning]
		Instead of hand-crafting task-specific features, representation learning trains neural network backbones to transform raw inputs $\mathbf{x}$ into continuous vector representations $\mathbf{z} \in \mathbb{R}^d$. Good representations disentangle underlying variation factors, compress noise, and make downstream tasks (classification, search, generation) trivial to solve using linear layers or distance metrics.
	\end{keyconcept}
	
	\textbf{Autoencoders and Latent Space Geometry.} An \textbf{Autoencoder} is a neural network trained to reconstruct its input through an architectural bottleneck:
	\begin{equation}
		\text{Encoder:}\quad \mathbf{z} = g_{\boldsymbol{\phi}}(\mathbf{x}) \in \mathbb{R}^d, \qquad \text{Decoder:}\quad \hat{\mathbf{x}} = f_{\boldsymbol{\theta}}(\mathbf{z}) \in \mathbb{R}^n .
	\end{equation}
	Minimizing reconstruction loss $\mathcal{L}_{\mathrm{rec}}(\mathbf{x}, \hat{\mathbf{x}}) = \|\mathbf{x} - \hat{\mathbf{x}}\|_2^2$ (or Binary Cross-Entropy) forces the bottleneck $\mathbf{z}$ to discard uninformative noise while preserving dominant structural variations:
	\begin{itemize}[nosep]
		\item \textbf{Undercomplete Autoencoders:} Enforcing bottleneck dimension $d < n$ prevents identity memorization. For linear encoders/decoders under MSE loss, the learned bottleneck subspace is equivalent to PCA.
		\item \textbf{Denoising Autoencoders (DAEs):} To prevent memorization when $d \ge n$, DAEs \parencite{vincent2008} corrupt inputs $\tilde{\mathbf{x}} \sim q(\tilde{\mathbf{x}} \mid \mathbf{x})$ via noise injection or random masking, training the decoder to reconstruct clean $\mathbf{x}$. DAEs learn a score-matching vector field $\nabla_{\mathbf{x}} \ln p(\mathbf{x})$ that projects off-manifold points back onto data manifold $\mathcal{M}$.
		\item \textbf{Variational Autoencoders (VAEs):} A probabilistic formulation \parencite{kingma2014} where the encoder outputs parameters $\boldsymbol{\mu}(\mathbf{x})$ and $\boldsymbol{\sigma}^2(\mathbf{x})$ of a latent Gaussian distribution $q_{\boldsymbol{\phi}}(\mathbf{z} \mid \mathbf{x}) = \mathcal{N}(\boldsymbol{\mu}, \mathrm{diag}(\boldsymbol{\sigma}^2))$. Training maximizes the \textbf{Evidence Lower Bound (ELBO)}:
		\begin{equation}
			\mathcal{L}_{\mathrm{ELBO}}(\boldsymbol{\phi}, \boldsymbol{\theta}; \mathbf{x}) = \mathbb{E}_{q_{\boldsymbol{\phi}}(\mathbf{z} \mid \mathbf{x})}\bigl[\ln p_{\boldsymbol{\theta}}(\mathbf{x} \mid \mathbf{z})\bigr] \;-\; D_{\mathrm{KL}}\!\left(q_{\boldsymbol{\phi}}(\mathbf{z} \mid \mathbf{x}) \,||\, p(\mathbf{z})\right),
		\end{equation}
		where $p(\mathbf{z}) = \mathcal{N}(\mathbf{0}, I)$ and $D_{\mathrm{KL}}$ penalizes latent distributions that deviate from standard normal priors. Reparameterization $\mathbf{z} = \boldsymbol{\mu} + \boldsymbol{\sigma} \odot \boldsymbol{\epsilon}$ ($\boldsymbol{\epsilon} \sim \mathcal{N}(\mathbf{0}, I)$) enables gradient backpropagation through stochastic sampling.
	\end{itemize}
	
	\textbf{Distributed Representations \& Word Embeddings.} Symbolic NLP represents words as orthogonal one-hot vectors $\mathbf{e}_i \in \{0, 1\}^{|V|}$, where dot products $\mathbf{e}_i^{\!\top}\mathbf{e}_j = 0$ for all $i \neq j$, failing to capture semantic similarity. \textbf{Distributed representations} \parencite{mikolov2013} embed discrete tokens into dense continuous spaces $\mathbb{R}^d$ ($d \ll |V|$), mapping semantic relationships directly to vector angles and distances.
	
	The \textbf{Word2Vec Skip-Gram} architecture predicts context tokens $w_{t+j}$ given target token $w_t$ within window size $c$. To avoid computing full-vocabulary softmax denominators over $|V|$, \textbf{Negative Sampling (SGNS)} reformulates training as binary logistic regressions distinguishing true context tokens from $K$ noise words sampled from unigram distribution $P_n(w)$:
	\begin{equation}
		\mathcal{L}_{\mathrm{SGNS}} = -\sum_{t=1}^{T} \sum_{-c \le j \le c, j \neq 0} \left[ \ln \sigma\!\left(\mathbf{v}'^{\!\top}_{w_{t+j}} \mathbf{v}_{w_t}\right) + \sum_{k=1}^{K} \mathbb{E}_{w_n \sim P_n(w)} \left[ \ln \sigma\!\left(-\mathbf{v}'^{\!\top}_{w_n} \mathbf{v}_{w_t}\right) \right] \right],
	\end{equation}
	where $\mathbf{v}_w$ and $\mathbf{v}'_w$ denote target and context vectors. An emergent property of SGNS spaces is \textbf{linear vector compositionality}: semantic analogies correspond to parallel displacement vectors ($\mathbf{v}_{\text{king}} - \mathbf{v}_{\text{man}} + \mathbf{v}_{\text{woman}} \approx \mathbf{v}_{\text{queen}}$).
	
	\textbf{Metric Learning \& Triplet Loss.} Representation learning can also optimize explicit distance metrics in embedding space. \textbf{Triplet Loss} trains encoders on anchor $\mathbf{a}$, positive $\mathbf{p}$ (same class), and negative $\mathbf{n}$ (different class) samples, enforcing that positive pairs remain closer than negative pairs by margin $\alpha > 0$:
	\begin{equation}
		\mathcal{L}_{\mathrm{triplet}} = \max\Bigl(0,\; \|\mathbf{z}_a - \mathbf{z}_p\|_2^2 - \|\mathbf{z}_a - \mathbf{z}_n\|_2^2 + \alpha\Bigr).
	\end{equation}

	\begin{practicebox}[Deep Autoencoder Feature Extraction \& Latent Space Analysis]
\begin{lstlisting}
import numpy as np
from sklearn.linear_model import LogisticRegression

# 1. Generate non-linear 10D synthetic dataset with 2D intrinsic manifold
rng = np.random.default_rng(42)
t = rng.uniform(-np.pi, np.pi, size=600)
# Non-linear 2D manifold embedded in 10D ambient space
X_raw = np.column_stack([np.sin(t), np.cos(t), t**2, np.sin(2*t), np.cos(2*t),
                         t, np.exp(t/5), np.sin(3*t), np.cos(3*t), t**3/10])
y_labels = (np.sin(t) + np.cos(2*t) > 0).astype(int)

# 2. Define 10D -> 4D -> 2D -> 4D -> 10D Deep Autoencoder in NumPy
W_e1 = rng.normal(0, np.sqrt(2/10), size=(10, 4)); b_e1 = np.zeros((1, 4))
W_e2 = rng.normal(0, np.sqrt(2/4), size=(4, 2));   b_e2 = np.zeros((1, 2))  # Bottleneck z in R^2
W_d1 = rng.normal(0, np.sqrt(2/2), size=(2, 4));   b_d1 = np.zeros((1, 4))
W_d2 = rng.normal(0, np.sqrt(2/4), size=(4, 10));  b_d2 = np.zeros((1, 10))

# 3. Train Autoencoder via SGD with MSE Reconstruction Loss
lr = 0.05
for epoch in range(400):
    # Forward Pass
    h1 = np.maximum(0, X_raw @ W_e1 + b_e1)           # Encoder L1 (ReLU)
    z  = h1 @ W_e2 + b_e2                             # Bottleneck representation z in R^2
    h2 = np.maximum(0, z @ W_d1 + b_d1)               # Decoder L1 (ReLU)
    X_hat = h2 @ W_d2 + b_d2                          # Reconstructed output in R^10
    
    # Backward Pass (Chain rule gradients)
    dX_hat = 2 * (X_hat - X_raw) / len(X_raw)
    dW_d2 = h2.T @ dX_hat;         db_d2 = dX_hat.sum(axis=0, keepdims=True)
    dh2   = (dX_hat @ W_d2.T) * (h2 > 0)
    dW_d1 = z.T @ dh2;            db_d1 = dh2.sum(axis=0, keepdims=True)
    dz    = dh2 @ W_d1.T
    dW_e2 = h1.T @ dz;             db_e2 = dz.sum(axis=0, keepdims=True)
    dh1   = (dz @ W_e2.T) * (h1 > 0)
    dW_e1 = X_raw.T @ dh1;        db_e1 = dh1.sum(axis=0, keepdims=True)
    
    # Parameter Updates
    W_e1 -= lr * dW_e1; b_e1 -= lr * db_e1; W_e2 -= lr * dW_e2; b_e2 -= lr * db_e2
    W_d1 -= lr * dW_d1; b_d1 -= lr * db_d1; W_d2 -= lr * dW_d2; b_d2 -= lr * db_d2

# 4. Extract 2D latent representation z and evaluate linear separability
z_latent = np.maximum(0, X_raw[:450] @ W_e1 + b_e1) @ W_e2 + b_e2
z_test   = np.maximum(0, X_raw[450:] @ W_e1 + b_e1) @ W_e2 + b_e2

clf_raw = LogisticRegression().fit(X_raw[:450], y_labels[:450])
clf_z   = LogisticRegression().fit(z_latent, y_labels[:450])

print(f"Epoch 400 Reconstruction MSE: {np.mean((X_hat - X_raw)**2):.4f}")
print(f"Linear Classifier Acc (Raw 10D Features):   {clf_raw.score(X_raw[450:], y_labels[450:]):.3f}")
print(f"Linear Classifier Acc (Learned 2D Latent z): {clf_z.score(z_test, y_labels[450:]):.3f}")
\end{lstlisting}
	\end{practicebox}
	
	\subsection{Architectures: CNNs, RNNs, and transformers}
	
	\textbf{Convolutional networks (CNNs)} exploit spatial structure in grid data (images, signals) via local receptive fields and weight sharing. A convolution slides a small learnable filter $K$ over input $I$:
	\begin{equation}
		(I \ast K)(i,j) = \sum_{m}\sum_{n} I(i{-}m,\, j{-}n)\,K(m,n).
	\end{equation}
	Sliding the filter confers built-in \textbf{translation equivariance} ($(I \ast K)(x - \Delta) = ((I \text{ shifted}) \ast K)(x)$): shifting an input shifts its feature map identically. Spatial pooling and hierarchical aggregation then build approximate \textbf{translation invariance} at higher layers. LeNet established document recognition \parencite{lecun1998}; AlexNet's GPU depth ignited modern computer vision \parencite{krizhevsky2012}.
	
	\textbf{Recurrent networks (RNNs)} handle sequences by carrying a hidden state:
	\begin{equation}
		\mathbf{h}_t = \tanh\!\bigl(W_h\,\mathbf{h}_{t-1} + W_x\,\mathbf{x}_t + \mathbf{b}\bigr).
	\end{equation}
	Training across long sequences suffers from vanishing gradients; the \textbf{LSTM}'s gated memory cells provided the decisive fix \parencite{hochreiter1997}, dominating sequential modeling for two decades.
	
	\textbf{Transformers} replaced sequential recurrent state updates with parallel \emph{self-attention}, positional encodings, and feed-forward transformations \parencite{vaswani2017}. Each token emits a query $Q$, a key $K$, and a value $V$; the output for each position is a weighted sum of all values, with weights given by query--key compatibility:
	\begin{equation}
		\mathrm{Attention}(Q, K, V) = \softmax\!\left(\frac{Q\,K^{\!\top}}{\sqrt{d_k}}\right) V .
	\end{equation}
	Because every token attends to all others in a single parallel step, transformers train vastly faster than RNNs on modern parallel hardware. However, computing pairwise attention scores across sequence length $n$ incurs $O(n^2)$ computational and memory complexity, creating a bottleneck for long \emph{context windows}.

	\begin{practicebox}[Scaled Dot-Product Attention from Scratch in NumPy]
\begin{lstlisting}
import numpy as np

def scaled_dot_product_attention(Q, K, V, mask=None):
    d_k = Q.shape[-1]
    scores = (Q @ K.T) / np.sqrt(d_k)
    if mask is not None:
        scores = np.where(mask == 0, -1e9, scores)
    weights = np.exp(scores - np.max(scores, axis=-1, keepdims=True))
    weights /= weights.sum(axis=-1, keepdims=True)  # Softmax
    return weights @ V, weights

# Demonstrate with 4 tokens of dimension 8
seq_len, d_k = 4, 8
Q = K = V = np.random.default_rng(0).standard_normal((seq_len, d_k))
causal_mask = np.tril(np.ones((seq_len, seq_len)))  # Lower-triangular mask
out, attn = scaled_dot_product_attention(Q, K, V, mask=causal_mask)
print("Attention Output Shape:", out.shape)
print("Causal Attention Weights Matrix:\n", np.round(attn, 3))
\end{lstlisting}
	\end{practicebox}

	\textbf{Multi-Head Attention (MHA).} Single-head attention averages relationships across the full representation space. To allow the network to jointly attend to information from different representation subspaces at different positions, \textbf{Multi-Head Attention} \parencite{vaswani2017} projects queries, keys, and values $h$ times with distinct learnable linear projections:
	\begin{equation}
		\mathrm{MultiHead}(Q, K, V) = \mathrm{Concat}(\mathrm{head}_1, \dots, \mathrm{head}_h) W^O,
	\end{equation}
	\begin{equation}
		\text{where}\quad \mathrm{head}_i = \mathrm{Attention}\bigl(Q W_i^Q,\; K W_i^K,\; V W_i^V\bigr),
	\end{equation}
	with projection parameter matrices $W_i^Q \in \mathbb{R}^{d_{\text{model}} \times d_k}$, $W_i^K \in \mathbb{R}^{d_{\text{model}} \times d_k}$, $W_i^V \in \mathbb{R}^{d_{\text{model}} \times d_v}$, and output projection $W^O \in \mathbb{R}^{h d_v \times d_{\text{model}}}$. In practice, $d_k = d_v = d_{\text{model}} / h$, keeping total computational cost comparable to single-head attention with full dimensionality.

	\begin{practicebox}[Multi-Head Attention from Scratch in NumPy]
\begin{lstlisting}
import numpy as np

def multi_head_attention(X, num_heads=4, d_model=64):
    B, N, _ = X.shape
    d_k = d_model // num_heads
    rng = np.random.default_rng(42)
    W_Q, W_K, W_V = [rng.normal(0, 0.1, size=(d_model, d_model)) for _ in range(3)]
    W_O = rng.normal(0, 0.1, size=(d_model, d_model))
    
    # Linear projection & reshape into (B, num_heads, N, d_k)
    Q = (X @ W_Q).reshape(B, N, num_heads, d_k).transpose(0, 2, 1, 3)
    K = (X @ W_K).reshape(B, N, num_heads, d_k).transpose(0, 2, 1, 3)
    V = (X @ W_V).reshape(B, N, num_heads, d_k).transpose(0, 2, 1, 3)
    
    # Parallel dot-product attention per head
    scores = (Q @ K.transpose(0, 1, 3, 2)) / np.sqrt(d_k)
    weights = np.exp(scores - np.max(scores, axis=-1, keepdims=True))
    weights /= weights.sum(axis=-1, keepdims=True)
    
    # Concatenate heads & final linear projection -> (B, N, d_model)
    heads_out = (weights @ V).transpose(0, 2, 1, 3).reshape(B, N, d_model)
    return heads_out @ W_O

X_in = np.random.default_rng(42).normal(size=(2, 8, 64))  # Batch 2, Seq 8, Dim 64
mha_out = multi_head_attention(X_in, num_heads=4, d_model=64)
print("Multi-Head Attention Output Shape:", mha_out.shape)  # (2, 8, 64)
\end{lstlisting}
	\end{practicebox}

	\textbf{Positional encodings.} Plain self-attention is permutation-invariant. To inject sequence order, \textbf{positional encodings} are added to input embeddings. The original Transformer used sinusoidal frequencies \parencite{vaswani2017}. Modern LLMs (e.g., Llama 3) predominantly use \textbf{Rotary Position Embeddings (RoPE)} \parencite{su2024}, rotating query and key vectors in complex 2D planes to encode relative distance directly into $Q K^\top$.

	\textbf{Residual connections, normalization, and optimization.} Deep Transformer stacks survive vanishing gradients by wrapping every sub-layer with a \textbf{residual (skip) connection} \parencite{he2016} and \textbf{Layer Normalization}. While original Transformers used \textbf{Post-Layer Normalization} ($\mathbf{y} = \mathrm{LayerNorm}(\mathbf{x} + \mathcal{F}(\mathbf{x}))$), modern foundation models standardly use \textbf{Pre-Layer Normalization} (Pre-LN):
	\begin{equation}
		\mathbf{y} = \mathbf{x} + \mathcal{F}\bigl(\mathrm{LayerNorm}(\mathbf{x})\bigr).
	\end{equation}
	Pre-LN preserves an unattenuated identity gradient highway ($\mathbf{x}$ passed directly to addition), enabling stable training across hundreds of layers without delicate warm-up schedules. Model optimization relies on \textbf{AdamW} \parencite{loshchilov2019} with learning rate warmup and cosine decay schedules.

	\begin{checkpoint}
		\begin{enumerate}
			\item Remove the nonlinearity from an MLP: what functions can the resulting network represent, and why?
			\item In an undercomplete Autoencoder with linear layers and MSE loss, what classical linear algebra technique does the bottleneck space replicate?
			\item Explain how Skip-Gram Negative Sampling (SGNS) overcomes the computational bottleneck of full-vocabulary softmax in Word2Vec.
			\item Why is Multi-Head Attention (MHA) preferred over single-head attention in modern Transformers?
		\end{enumerate}
	\end{checkpoint}

	\begin{sourcesbox}
		\small
		\begin{itemize}[nosep]
			\item \fullcite{rosenblatt1958}
			\item \fullcite{rumelhart1986}
			\item \fullcite{hornik1989}
			\item \fullcite{nair2010}
			\item \fullcite{bengio2013}
			\item \fullcite{vincent2008}
			\item \fullcite{lecun1998}
			\item \fullcite{hochreiter1997}
			\item \fullcite{he2016}
			\item \fullcite{vaswani2017}
			\item \fullcite{glorot2010}
			\item \fullcite{he2015}
		\end{itemize}
	\end{sourcesbox}

	% =====================================================================
	\clearpage
	\section{Foundation Models and Generative AI}
	% =====================================================================
	
	\begin{modulecard}{\faRobot\ Module 5 --- at a glance \hfill \small\normalfont (Estimated Study: 8 hours \textbullet{} Level: Advanced)}{primary!80!black}
		\textbf{Focus.} The self-supervised foundation model paradigm: subword tokenization, SSL pretext objectives, scaling laws, generative models (diffusion, LLMs), fine-tuning alignment (LoRA, RLHF, DPO), high-performance inference (KV-caching, RAG), and responsible AI revisited.
		
		\medskip
		\textbf{Key concepts.} Tokenization (BPE), self-supervised learning, InfoNCE loss, Causal LM, Masked LM, scaling laws, diffusion models, SFT, LoRA/QLoRA, RLHF, DPO, KV-caching, RAG, and alignment failures.
		
		\medskip
		\textbf{After this module you will be able to} explain how raw text is converted into tokens, derive contrastive InfoNCE and DPO loss formulations, explain compute-optimal scaling laws, implement parameter-efficient fine-tuning via LoRA, trace LLM next-token generation step-by-step, and audit foundation models against ethical safety dimensions.
	\end{modulecard}
	
	\subsection{Self-supervised learning and subword tokenization}
	
	The transition from classical machine learning to modern foundation models rests fundamentally on \textbf{Self-Supervised Learning (SSL)}. Classical supervised learning achieves high accuracy but is strictly bottlenecked by expensive, non-scalable human annotations $(\mathbf{x}_i, y_i)$. Traditional unsupervised learning (K-Means, PCA) scales to unlabelled data but struggles to discover task-relevant high-level semantics without targets. Self-supervised learning bridges this divide by turning unlabelled raw data into its own supervisor: the model automatically constructs artificial target signals (\emph{pretext tasks}) directly from intrinsic structural patterns within raw inputs $\mathbf{x}$.

	\textbf{Tokenization: Bridging Raw Text and Model Inputs.} Foundation models do not process raw string characters or whole dictionary words directly. Whole-word vocabularies suffer from out-of-vocabulary (OOV) failures and massive matrix dimensions, while character-level models yield long sequence lengths. Modern LLMs use \textbf{Subword Tokenization} --- primarily \textbf{Byte-Pair Encoding (BPE)} --- which iteratively merges the most frequent adjacent character or byte pairs into subwords.

	\begin{practicebox}[Byte-Pair Encoding (BPE) Tokenization Step from Scratch]
\begin{lstlisting}
from collections import defaultdict

def get_stats(vocab):
    pairs = defaultdict(int)
    for word, freq in vocab.items():
        symbols = word.split()
        for i in range(len(symbols) - 1):
            pairs[symbols[i], symbols[i+1]] += freq
    return pairs

def merge_vocab(pair, v_in):
    v_out = {}
    bigram = ' '.join(pair)
    replacement = ''.join(pair)
    for word in v_in:
        w_out = word.replace(bigram, replacement)
        v_out[w_out] = v_in[word]
    return v_out

# Synthetic toy vocabulary with word frequencies
vocab = {'l o w </w>': 5, 'l o w e r </w>': 2, 'n e w e s t </w>': 6, 'w i d e s t </w>': 3}
print("Initial Subword Vocab:", vocab)
for step in range(3):
    pairs = get_stats(vocab)
    best_pair = max(pairs, key=pairs.get)
    vocab = merge_vocab(best_pair, vocab)
    print(f"Step {step+1} Best Pair {best_pair} -> Merged Vocab: {vocab}")
\end{lstlisting}
	\end{practicebox}

	\begin{keyconcept}[Self-Supervised Learning: The Foundation Model Engine]
		Self-Supervised Learning converts unlabelled raw data (text, images, audio) into supervised learning targets by hiding, masking, or augmenting parts of the input and training a network to predict the missing signal. SSL provides an unlimited training signal, turning web-scale raw datasets into universal pre-trained representations.
	\end{keyconcept}
	
	\textbf{Taxonomy of SSL Pretext Objectives.} SSL transforms raw signals into target supervision through three primary mathematical formulations:
	\begin{enumerate}
		\item \textbf{Causal (Autoregressive) Language Modeling (CLM):} The model predicts token $x_t$ conditioned on preceding historical context $\mathbf{x}_{1:t-1}$, minimizing negative log-likelihood over sequence length $T$:
		\begin{equation}
			\mathcal{L}_{\mathrm{CLM}}(\boldsymbol{\theta}) = -\sum_{t=1}^{T} \ln P_{\boldsymbol{\theta}}(x_t \mid x_1, x_2, \dots, x_{t-1}).
		\end{equation}
		Pioneered at scale by the GPT family \parencite{radford2018, brown2020}, CLM creates generative autoregressive models capable of next-token prediction and zero/few-shot in-context learning.
		
		\item \textbf{Masked Reconstruction (MLM \& MAE):} A subset of positions $\mathcal{M}$ (typically 15\% in text, 75\% in vision) is corrupted or masked out. The network predicts original values $x_i$ using bidirectional context $\mathbf{x}_{\setminus \mathcal{M}}$:
		\begin{equation}
			\mathcal{L}_{\mathrm{MLM}}(\boldsymbol{\theta}) = -\sum_{i \in \mathcal{M}} \ln P_{\boldsymbol{\theta}}(x_i \mid \mathbf{x}_{\setminus \mathcal{M}}).
		\end{equation}
		Pioneered by BERT \parencite{devlin2019} in NLP and Masked Autoencoders (MAE) \parencite{he2022} in computer vision, masked reconstruction forces bidirectional attention layers to model deep contextual dependencies.
		
		\item \textbf{Contrastive Learning \& InfoNCE Loss:} Beyond text, SSL for vision and multimodal data generates two stochastic augmented views $(\tilde{\mathbf{x}}_i, \tilde{\mathbf{x}}_i')$ of the same image (positive pair) via cropping, color jittering, and blurring. Encoders project views into normalized latent vectors $\mathbf{z}_i = g(f(\tilde{\mathbf{x}}_i))$. The \textbf{InfoNCE Loss} \parencite{oord2018, chen2020} pulls positive pairs together while pushing negative views $\mathbf{z}_k^-$ (other images in the batch) apart:
		\begin{equation}
			\mathcal{L}_{\mathrm{InfoNCE}} = -\sum_{i=1}^{N} \ln \frac{\exp\left(\mathrm{sim}(\mathbf{z}_{i}, \mathbf{z}_{i}^+) / \tau\right)}{\exp\left(\mathrm{sim}(\mathbf{z}_{i}, \mathbf{z}_{i}^+) / \tau\right) + \sum_{k \neq i} \exp\left(\mathrm{sim}(\mathbf{z}_{i}, \mathbf{z}_{k}^-) / \tau\right)},
		\end{equation}
		where $\mathrm{sim}(\mathbf{u}, \mathbf{v}) = \frac{\mathbf{u}^{\!\top}\mathbf{v}}{\|\mathbf{u}\|_2 \|\mathbf{v}\|_2}$ is cosine similarity and temperature hyperparameter $\tau > 0$ controls the penalty hardness assigned to difficult negative samples.
	\end{enumerate}

	\textbf{Non-contrastive SSL \& Collapse Prevention.} Contrastive learning requires large batch sizes (e.g., 4096 in SimCLR \parencite{chen2020}) to supply sufficient negative samples. Non-contrastive methods (BYOL, SimSiam) eliminate negative pairs entirely, training an online network to predict target network representations. Representation collapse (where the encoder outputs a constant vector $\mathbf{z}_i = \mathbf{c}$) is prevented using asymmetrical predictor heads, moving-average target networks, and \emph{stop-gradient} operations ($\mathrm{sg}[\cdot]$).
	
	\textbf{The Symbiotic Relationship between SSL and Transformers.} The simultaneous rise of Transformers \parencite{vaswani2017} and SSL represents a profound technological synergy. SSL provided an infinite supply of training supervision from web text, while the Transformer provided the exact architectural mechanisms required to ingest that scale:
	\begin{itemize}
		\item \textbf{Elimination of Recurrent Bottlenecks:} RNNs process input sequences sequentially ($\mathbf{h}_t = f(\mathbf{h}_{t-1}, \mathbf{x}_t)$), requiring $O(N)$ sequential operations during training. Backpropagation Through Time (BPTT) creates severe memory bottlenecks. Transformers compute all pairwise self-attention interactions in parallel ($O(1)$ sequential depth), unlocking full GPU throughput across trillions of SSL tokens.
		\item \textbf{Unrestricted Inductive Bias \& Scaling Laws:} CNNs enforce local spatial kernels, and RNNs enforce strict temporal decay. While helpful on small datasets, rigid inductive biases cap model capacity on massive datasets. Transformers replace rigid priors with dynamic soft self-attention. When combined with SSL, Transformer capacity scales predictably according to empirical power laws without early performance saturation \parencite{kaplan2020, hoffmann2022}.
		\item \textbf{The Pretrain--Finetune Paradigm Shift:} SSL converts raw unlabelled data into universal cognitive representations. The modern AI paradigm uses SSL to pretrain massive base Transformer models, which transfer efficiently to downstream tasks via Supervised Fine-Tuning (SFT), Parameter-Efficient Adaptation (LoRA \parencite{hu2021}), or zero/few-shot in-context learning \parencite{brown2020}.
	\end{itemize}

	\begin{practicebox}[SimCLR Contrastive Learning with InfoNCE Loss from Scratch]
\begin{lstlisting}
import numpy as np

# 1. Define vectorized InfoNCE Loss with Temperature Scaling (SimCLR style)
def infonce_loss(z1, z2, temperature=0.1):
    """
    z1, z2: (N, d) L2-normalized embedding projections for 2 augmented views
    """
    N = len(z1)
    # Concatenate both views -> 2N x d matrix
    z = np.vstack([z1, z2])
    
    # Compute 2N x 2N cosine similarity matrix (since vectors are L2-normalized)
    sim_matrix = (z @ z.T) / temperature
    
    # Mask out self-similarity along main diagonal
    np.fill_diagonal(sim_matrix, -1e9)
    
    # Positive pair target indices: row i pairs with (i + N) % (2N)
    labels = np.arange(2 * N)
    labels[:N] += N
    labels[N:] -= N
    
    # Numerically stable Log-Softmax Cross-Entropy Loss
    max_sim = np.max(sim_matrix, axis=1, keepdims=True)
    exp_sim = np.exp(sim_matrix - max_sim)
    log_prob = sim_matrix - max_sim - np.log(np.sum(exp_sim, axis=1, keepdims=True))
    
    loss = -np.mean(log_prob[np.arange(2 * N), labels])
    return loss

# 2. Simulate 2 augmented views of 4 samples in 8-dimensional projection space
rng = np.random.default_rng(42)
N, d = 4, 8
# Base representations
z_base = rng.normal(size=(N, d))
z_base /= np.linalg.norm(z_base, axis=1, keepdims=True)

# Positive views (small random augmentation perturbation)
z1 = z_base + 0.05 * rng.normal(size=(N, d))
z2 = z_base + 0.05 * rng.normal(size=(N, d))

# L2 normalize projection vectors
z1 /= np.linalg.norm(z1, axis=1, keepdims=True)
z2 /= np.linalg.norm(z2, axis=1, keepdims=True)

loss_aligned = infonce_loss(z1, z2, temperature=0.1)

# Unaligned random representations baseline
z_rand = rng.normal(size=(N, d))
z_rand /= np.linalg.norm(z_rand, axis=1, keepdims=True)
loss_unaligned = infonce_loss(z1, z_rand, temperature=0.1)

print(f"InfoNCE Loss (Aligned Positive Pairs):    {loss_aligned:.4f}")
print(f"InfoNCE Loss (Random Unaligned Pairs):   {loss_unaligned:.4f}")
\end{lstlisting}
	\end{practicebox}
	
	\subsection{Scaling laws and generative architectures}

	\textbf{Latent Space \& Diffusion Models.} Generative models compress data into lower-dimensional latent representations. Variational Autoencoders \parencite{kingma2014} and GANs \parencite{goodfellow2014} pioneered latent generation. \textbf{Diffusion Models} \parencite{ho2020} dominate image/audio synthesis by training U-Nets or Vision Transformers to reverse a forward Gaussian noise process $q(\mathbf{x}_t \mid \mathbf{x}_{t-1})$.
	
	\textbf{Large Language Models \& Causal Masking.} Autoregressive LLMs predict the next token via cross-entropy loss over vocabulary $V$:
	\begin{equation}
		p(y = j \mid \mathbf{x}) = \frac{\exp(z_j)}{\sum_{k=1}^{V} \exp(z_k)}, \qquad \mathcal{L}_{\text{CE}} = -\log p(y_{\text{true}} \mid \mathbf{x}).
	\end{equation}
	Causal attention masking ($M_{ij} = -\infty$ for $j > i$) prevents tokens from viewing future positions.
	
	\textbf{Compute-Optimal Scaling Laws.} Pretraining consumes trillions of tokens. Empirical \textbf{scaling laws} \parencite{kaplan2020} showed power-law scaling of performance with compute, parameters, and tokens. Compute-optimal \textbf{Chinchilla scaling analysis} \parencite{hoffmann2022} demonstrated that model capacity and training token count should scale in roughly equal $1:1$ proportion under compute-optimal budgets.

	\subsection{Fine-tuning, parameter efficiency, and preference alignment}

	Following pretraining, alignment transforms raw base models into helpful assistants:
	\begin{enumerate}[nosep]
		\item \textbf{Supervised Fine-Tuning (SFT):} Fine-tuning on curated demonstration pairs.
		\item \textbf{Preference Optimization (RLHF \& DPO):} Reinforcement Learning from Human Feedback \parencite{christiano2017, ouyang2022} or Direct Preference Optimization (DPO) \parencite{rafailov2023} optimizes policy preference loss directly without fitting an explicit reward model.
		\item \textbf{Parameter-Efficient Fine-Tuning (PEFT / LoRA):} \textbf{LoRA} \parencite{hu2021} freezes pretrained weight matrices $W_0 \in \mathbb{R}^{d \times k}$ and injects trainable low-rank decomposition matrices $\Delta W = B A$ ($B \in \mathbb{R}^{d \times r}, A \in \mathbb{R}^{r \times k}, r \ll \min(d,k)$), reducing trainable parameters by $>99\%$. \textbf{QLoRA} \parencite{dettmers2024} quantizes base model weights to 4-bit NormalFloat while preserving full-precision low-rank adapters.
	\end{enumerate}

	\subsection{Foundation model inference, KV-caching, and RAG}

	At inference, token sampling controls decoding:
	\begin{itemize}[nosep]
		\item \textbf{Temperature Scaling \& Greedy Decoding:} Softmax $\mathrm{softmax}(\mathbf{z}/T)$ controls distribution entropy. As $T \to 0^+$, the distribution converges to a one-hot vector at $\arg\max \mathbf{z}$, which code libraries alias as \texttt{temperature=0} (greedy decoding).
		\item \textbf{Sampling Strategies:} Top-$k$ truncates sampling to the $k$ most probable tokens; Top-$p$ (nucleus) samples from the smallest candidate set comprising cumulative probability $p$.
		\item \textbf{Inference Acceleration:} High-throughput inference relies on \textbf{KV Caching}\footnote{KV Caching stores past $K$ and $V$ tensors in memory, making incremental attention computation per generated token linear in context length $O(N)$ instead of $O(N^2)$.}, continuous batching, quantization (INT8/INT4), and speculative decoding.
		\item \textbf{Retrieval-Augmented Generation (RAG):} RAG \parencite{lewis2020} grounds generation by retrieving relevant document snippets from vector databases into prompt context.
	\end{itemize}

	\begin{keyconcept}[How an LLM writes one token]
		\begin{enumerate}
			\item Tokenise the prompt into sub-word tokens.
			\item Embed each token into a vector (plus RoPE positional encoding).
			\item Pass the sequence through stacked Pre-LN transformer blocks (causal self-attention + FFN).
			\item Map final position hidden state through linear head and softmax over vocabulary $V$.
			\item Sample next token via Temperature / Top-$p$ or select argmax greedy token ($T \to 0^+$).
			\item Append token to sequence and update KV-cache for the next step.
		\end{enumerate}
	\end{keyconcept}

	\subsection{Ethics, alignment, and safety revisited: the generative AI frontier}

	Module~1 introduced six baseline dimensions of AI ethics: algorithmic bias, privacy, alignment, copyright/provenance, environmental footprint, and automation bias. Generative AI and foundation models exacerbate these challenges while introducing distinct novel failure modes:

	\begin{ethicsbox}[Ethics \& Safety Spotlight: Foundation Model Risk Checklist]
		\begin{itemize}[nosep]
			\item \textbf{Hallucinations \& Grounding Failure:} Fluent syntax masks ungrounded factual inaccuracies. RLHF can exacerbate "sycophancy" where models echo user biases rather than truth.
			\item \textbf{Data Provenance \& Copyright:} Web-scale pretraining relies on copyrighted literary and visual art without attribution or compensation, driving legal and policy battles.
			\item \textbf{Watermarking \& Synthetic Content Identification:} As AI-generated text, image, and deepfake video proliferate, cryptographic provenance and statistical watermarking become vital for info-ecosystem integrity.
			\item \textbf{Environmental \& Compute Externalities:} Training modern foundation models consumes gigawatt-hours of electricity and significant cooling water \parencite{strubell2019, patterson2021}, highlighting the need for efficient architectures and compute-optimal training.
		\end{itemize}
	\end{ethicsbox}

	\begin{checkpoint}
		\begin{enumerate}
			\item Why does Self-Supervised Learning (SSL) depend so critically on the parallel self-attention mechanism of Transformers compared to sequential RNNs?
			\item Mathematically, what role does temperature $\tau$ play in the InfoNCE loss for contrastive learning?
			\item What key trade-off does subword tokenization (e.g. Byte-Pair Encoding) balance compared to whole-word or character-level tokenization?
			\item According to Chinchilla scaling laws, what is the compute-optimal ratio between scaling model parameters vs.\ scaling training tokens?
			\item Explain how LoRA achieves $>99\%$ parameter reduction during fine-tuning via low-rank matrix decomposition.
		\end{enumerate}
	\end{checkpoint}

	\begin{sourcesbox}
		\small
		\begin{itemize}[nosep]
			\item \fullcite{vaswani2017}
			\item \fullcite{oord2018}
			\item \fullcite{chen2020}
			\item \fullcite{devlin2019}
			\item \fullcite{radford2018}
			\item \fullcite{dosovitskiy2021}
			\item \fullcite{he2022}
			\item \fullcite{su2024}
			\item \fullcite{mikolov2013}
			\item \fullcite{kingma2014}
			\item \fullcite{goodfellow2014}
			\item \fullcite{ho2020}
			\item \fullcite{brown2020}
			\item \fullcite{hoffmann2022}
			\item \fullcite{ouyang2022}
			\item \fullcite{rafailov2023}
			\item \fullcite{hu2021}
			\item \fullcite{dettmers2024}
			\item \fullcite{lewis2020}
			\item \fullcite{strubell2019}
			\item \fullcite{patterson2021}
			\item \fullcite{bender2021}
		\end{itemize}
	\end{sourcesbox}

	% =====================================================================
	\clearpage
	\section*{Where to go next}
	\markright{Where to go next}
	\phantomsection
	\addcontentsline{toc}{section}{Where to go next}
	% =====================================================================
	
	The guide ends; the field does not. A few well-marked paths, matched to what you now know:
	\begin{itemize}
		\item \textbf{Strengthen the mathematics:} \parencite{strang2019, wasserman2004, boyd2004}.
		\item \textbf{Go deep on classical ML:} \parencite{hastie2009, bishop2006, pedregosa2011}.
		\item \textbf{Deep learning proper:} \parencite{goodfellow2016, bishop2024, zhang2023}.
		\item \textbf{Build with LLMs, LoRA \& RAG:} Hugging Face ecosystem \parencite{wolf2020}, LoRA \parencite{hu2021}, and RAG \parencite{lewis2020}, read alongside \parencite{vaswani2017, brown2020, ouyang2022, rafailov2023}.
		\item \textbf{Stay grounded:} \parencite{christian2020, bender2021}.
	\end{itemize}

	% =====================================================================
	\clearpage
	\section*{Practical Coding Projects Series}
	\markright{Practical Coding Projects Series}
	\phantomsection
	\addcontentsline{toc}{section}{Practical Coding Projects Series}
	% =====================================================================
	
	To accompany the theoretical foundation developed across the five modules, a complete series of six substantial hands-on coding projects is provided in the repository's \texttt{projects/} directory (\url{https://github.com/DanielBerns/AI-course}). Each project bridges theoretical concepts to runnable Python implementation using industry-standard frameworks (\texttt{NumPy}, \texttt{Pandas}, \texttt{Scikit-Learn}, \texttt{XGBoost}, \texttt{PyTorch}, and \texttt{Transformers}), with explicit emphasis on data cleaning, exploratory visualization, model diagnostics, resumable training, and inference pipelines.

	% ---------------------------------------------------------------------
	\subsection*{Project 1: Data Pipeline, Exploratory Analysis \& Algorithmic Bias Audit}
	% ---------------------------------------------------------------------
	\textbf{Curriculum Mapping:} Module 1 (The AI Landscape, Compute/Cloud, Data Pipelines, Ethics \& Safety).

	\textbf{Overview \& Learning Objectives:} Hands-on data engineering and ethical auditing pipeline preceding model training on clinical dataset.

	% ---------------------------------------------------------------------
	\subsection*{Project 2: Math Engine from Scratch (Linear Algebra, Stats \& Optimization)}
	% ---------------------------------------------------------------------
	\textbf{Curriculum Mapping:} Module 2 (Mathematics and Data Foundations). Pure NumPy matrix calculus, PCA, Adam/AdamW optimizers.

	% ---------------------------------------------------------------------
	\subsection*{Project 3: Classical Machine Learning Pipeline \& Diagnostics}
	% ---------------------------------------------------------------------
	\textbf{Curriculum Mapping:} Module 3 (Classical Machine Learning). Supervised classification, DBSCAN/K-Means clustering, IsoForest anomaly detection, learning curves.

	% ---------------------------------------------------------------------
	\subsection*{Project 4: Deep Learning Foundations \& Resumable Training Pipeline}
	% ---------------------------------------------------------------------
	\textbf{Curriculum Mapping:} Module 4 (Neural Networks, Backpropagation, PyTorch/Keras). Manual matrix backprop engine, PyTorch MLP, checkpoint manager.

	% ---------------------------------------------------------------------
	\subsection*{Project 5: Computer Vision \& Sequential Architectures (CNNs \& LSTMs)}
	% ---------------------------------------------------------------------
	\textbf{Curriculum Mapping:} Module 4 (CNNs, RNNs/LSTMs, Multi-Head Attention). PyTorch 2D CNN with Pre-LN Residual Skip Connections and PyTorch LSTM.

	% ---------------------------------------------------------------------
	\subsection*{Project 6: Transformer Architectures \& LLM Fine-Tuning Pipeline}
	% ---------------------------------------------------------------------
	\textbf{Curriculum Mapping:} Module 5 (Foundation Models, Generative AI \& LLMs). Causal self-attention from scratch, RoPE, Hugging Face fine-tuning (SFT/LoRA), next-token sampling engine.

	% =====================================================================
	\clearpage
	\section*{Appendix: Answer Key for Check Your Understanding}
	\markright{Appendix: Answer Key}
	\phantomsection
	\addcontentsline{toc}{section}{Appendix: Answer Key for Check Your Understanding}
	% =====================================================================

	This appendix provides concise model answers for all self-check exercises across the five modules.

	\subsection*{Module 1: Introduction and the AI Landscape}
	\begin{enumerate}
		\item \textbf{Classification:} Chess engine, spam filter, and autonomous reasoning demos are all Narrow AI.
		\item \textbf{GPU vs CPU throughput:} GPUs possess thousands of concurrent lightweight arithmetic cores optimized for parallel matrix tensor operations ($A \mathbf{x} + \mathbf{b}$).
		\item \textbf{Job screening harms \& mitigations:} (1) Historical bias propagation $\to$ audit for Equalized Odds / proxy features; (2) Representation gaps $\to$ human oversight \& threshold calibration.
	\end{enumerate}

	\subsection*{Module 2: Mathematics and Data Foundations}
	\begin{enumerate}
		\item \textbf{Eigenvector of Covariance Matrix:} Specifies principal direction of maximum variance projection without linear distortion.
		\item \textbf{Bayes' Theorem with 10\% Prevalence:} $P(\text{disease}\mid +) = \frac{0.99 \times 0.10}{(0.99 \times 0.10) + (0.05 \times 0.90)} \approx 0.6875$ (68.8\%).
		\item \textbf{Divergence with $\eta=1.1$ for $J(w)=w^2$:} Step $w_{t+1} = -1.2 w_t$ flips sign and expands magnitude by $1.2\times$, diverging to $\pm\infty$.
	\end{enumerate}

	\subsection*{Module 3: Classical Machine Learning}
	\begin{enumerate}
		\item \textbf{Cancer Screening Metric:} Prioritize Recall to minimize False Negatives; precision drops (more false alarms requiring follow-up tests).
		\item \textbf{K-Means on Concentric Rings:} K-Means assumes convex spherical clusters; concentric rings require density-based clustering (DBSCAN) or spectral methods.
		\item \textbf{K-Armed Bandit $\epsilon$ Trade-off:} Higher $\epsilon$ accelerates early exploration but caps asymptotic long-term return due to persistent suboptimal random actions.
		\item \textbf{Cross-Validation Score Variance (0.62 to 0.91):} Indicates model instability or fold sample imbalance; use stratified $k$-fold, stronger regularization, or more data.
	\end{enumerate}

	\subsection*{Module 4: Neural Networks and Representation Learning}
	\begin{enumerate}
		\item \textbf{MLP without Nonlinearity:} Collapses mathematically into a single linear map $\mathbf{a}^{(L)} = W_{\text{effective}}\mathbf{x}$, restricted to linear hyperplanes.
		\item \textbf{Linear Undercomplete Autoencoder Subspace:} Under MSE loss, a linear undercomplete autoencoder bottleneck learns the exact principal subspace as Principal Component Analysis (PCA).
		\item \textbf{Skip-Gram Negative Sampling (SGNS):} Converts expensive full-vocabulary softmax over $|V|$ into $K+1$ binary logistic regressions per step, replacing $O(|V|)$ denominator summations with $O(K)$ binary updates.
		\item \textbf{Multi-Head Attention Advantage:} Projects inputs into $h$ distinct subspace heads in parallel, enabling the model to simultaneously attend to different aspects (e.g. syntactic agreement vs semantic reference) at different positions.
	\end{enumerate}

	\subsection*{Module 5: Foundation Models and Generative AI}
	\begin{enumerate}
		\item \textbf{SSL \& Transformer Synergy:} SSL relies on pseudo-tasks (CLM/MLM) over unlabelled web-scale corpora containing trillions of tokens. Sequential RNNs process inputs sequentially ($O(N)$ steps) with memory bottlenecks in BPTT, making web-scale training computationally intractable. Transformers compute all pairwise token interactions in parallel ($O(1)$ sequential operations) and lack rigid local/temporal inductive biases, enabling parameter capacity to absorb massive scale without early performance saturation.
		\item \textbf{InfoNCE Temperature $\tau$:} Scales cosine similarities before log-softmax, controlling the relative gradient penalty weight assigned to hard negative samples vs.\ easy negatives.
		\item \textbf{Subword Tokenization Trade-off:} Balances vocabulary size and sequence length by keeping frequent words intact while breaking rare words into subword pieces, avoiding out-of-vocabulary (OOV) tokens.
		\item \textbf{Chinchilla Scaling Ratio:} Specifies that for optimal compute allocation, parameter count and training token count should scale at an equal $1:1$ ratio.
		\item \textbf{LoRA Parameter Reduction:} Replaces full weight matrix update $\Delta W \in \mathbb{R}^{d \times k}$ with low-rank decomposition $B A$ ($B \in \mathbb{R}^{d \times r}, A \in \mathbb{R}^{r \times k}$ with $r \ll \min(d, k)$), drastically reducing trainable parameters from $d \cdot k$ to $r(d + k)$.
	\end{enumerate}

	% =====================================================================
	\clearpage
	\section*{Glossary}
	\markright{Glossary}
	\phantomsection
	\addcontentsline{toc}{section}{Glossary}
	% =====================================================================
	
	\begin{table}[ht]
		\centering
		\small\sffamily
		\renewcommand{\arraystretch}{1.25}
		\rowcolors{2}{black!2}{black!5}
		\caption{Quick reference for the vocabulary used in this guide.}
		\begin{tabular}{@{}L{3.8cm} L{10.8cm}@{}}
			\toprule
			\rowcolor{primary}
			\textcolor{white}{\textbf{Term}} & \textcolor{white}{\textbf{Meaning}} \\
			\midrule
			Artificial intelligence (AI) & Broad field: systems that perform tasks associated with intelligent behaviour. \\
			Machine learning (ML) & Sub-field: algorithms whose performance on a task improves from data. \\
			Deep learning (DL) & ML with multi-layer neural networks that learn hierarchical representations. \\
			Feature / label & Input measurement / target answer supplied during supervised training. \\
			Model & Parameterised function fitted to data. \\
			Loss function & Scalar score of single-example error; cost function averages loss over data. \\
			Gradient & Vector of partial derivatives indicating the direction of steepest ascent. \\
			Backpropagation & Efficient chain-rule algorithm computing loss gradients for all weights layer by layer. \\
			Epoch & One complete pass over the entire training dataset during model optimization. \\
			Hyperparameter & Configuration variable set before training rather than learned from data. \\
			Regularization & Techniques ($\ell_1/\ell_2$, dropout, weight decay) constraining model complexity to prevent overfitting. \\
			Vanishing gradient & Exponential decay of backpropagated error gradients toward zero in deep networks. \\
			Attention & Mechanism dynamically weighting sequence representations based on query relevance. \\
			Multi-Head Attention & Parallel attention mechanisms operating over distinct linear projection subspaces. \\
			Inductive bias & Set of assumptions an algorithm uses to predict outputs for unseen inputs. \\
			Overfitting & Memorising training data; high training performance, poor generalization. \\
			Embedding & Dense vector representation encoding semantic relationships in continuous space. \\
			Tokenization & Segmenting raw text into discrete integer vocabulary tokens (e.g. via BPE). \\
			Token & Sub-word unit that language models process and generate. \\
			Parameter & Learnable weight/bias inside a model; LLMs contain billions of parameters. \\
			Fine-tuning & Adapting a pretrained model to a specific task via SFT, LoRA, or DPO. \\
			Hallucination & Model output that is fluent and syntactically convincing but factually false. \\
			\bottomrule
		\end{tabular}
	\end{table}
	
	% =====================================================================
	\clearpage
	\printbibliography[title={References}, heading=bibintoc]
	% =====================================================================
	
	\vfill
	\begin{center}
		{\small\sffamily\color{textgray}\rule{0.3\linewidth}{0.6pt}\\[0.5em]
			--- End of document ---}
	\end{center}
	
\end{document}